% SOURCE docs/release/2026-10-29/source/docs/causal_banking_release_candidate_2026-10-29.tex CHAR 27043
\chapter{Comparing assistance with credible alternatives}
An assistant should be judged against the work an analyst could already
do. Evaluation therefore needs simple baselines, tasks with defensible answers,
and a record of unresolved cases and correction costs. The following studies
ask what benchmarks, retrieval, code generation and simulated populations can
establish about that comparison.



% SOURCE docs/release/2026-10-29/source/docs/writing-repair/2026-10-29/src/8d4e1eb50e94271d.tex CHAR 0
\section{Constructed causal labels, semantic refusal and the unit of evaluation}
\label{exec17:benchmark-development}

The preceding chapter established what an evaluation must preserve. This chapter asks which simpler or more established method supplies a credible comparison at each stage. Constructed labels and literature retrieval test the reference; coding and repair studies test calculation; factor extraction and simulated respondents test whether the inputs represent the intended population. The final program-of-thought study returns to the full chain, where saved accuracy, scientific meaning and executable reuse can be assessed separately.

A causal assistant can improve the wording of an answer while leaving its
decision content unchanged. It can also repeat a familiar causal story more
fluently than an analyst without establishing that the story applies to the
bank's observations. Evaluation therefore needs to distinguish recognizing a
stipulated mechanism, identifying an effect from supplied information and
choosing a useful action. Two further benchmarks make this distinction
concrete: the published CELLO study by Chen and colleagues and the first
Research Square version of CausalReasonBench by Cherif. The former uses visual
scenes and constructed graphs; the latter uses textual scenarios and automatic
reasoning scores. Their results concern their respective tasks and historical
model settings. Neither is a measurement of a causal banking
workflow.~\citep{exec17_benchmark_cello,exec17_benchmark_cherif}

Consider the proposed deposit-pricing assistant. A manager asks whether a
retention offer will prevent withdrawals. A useful response must distinguish
the association between receiving an offer and staying with the bank from the
effect of offering the same terms to an eligible customer. It must also state
what the information cannot establish. A benchmark that supplies a complete
causal mechanism tests whether the assistant follows that mechanism. A
benchmark that supplies only an observational vignette may instead need an
answer that preserves several possible mechanisms. Those are different
questions even when the surface wording and requested answer are identical.
Constructed labels can be valuable for the first question; their construction
does not resolve the second.

\subsection{What a scene-derived causal answer establishes}

CELLO builds causal reasoning questions from Visual Genome relations. The
construction matches subgraphs to graph templates and uses language-model
generation to produce questions and distractors. Its prompting procedure
separates entity extraction, graph construction, task recognition and causal
knowledge. This organization is relevant to the proposal because a banking
assistant also needs to turn entities and relationships into an explicit
question before asking a numerical engine to answer it. The distinction is
that a relation visible in an image, or recorded in a customer document, need
not determine the intervention semantics assigned by a template. Static
co-occurrence, a plausible physical explanation and a warranted counterfactual
are distinct kinds of information. The benchmark's annotation and image
limitations therefore constrain the interpretation of its scores, rather than
being mere implementation details.~\citep{exec17_benchmark_cello}

The paper reports a test set of 1,200 questions, with 100 selected for each of
twelve tasks. Its Appendix B.2 describes a separate validity check of 120
instances, ten per task, by two annotators: 91.7\% were judged valid, with
reported Cohen agreement of 84.1\%. These figures describe a question-validity
review. They do not turn each template's intervention or counterfactual claim
into experimentally established truth. For banking, the analogous review
would check that the prompt faithfully expresses an agreed mechanism and
information set. It would still leave open whether that mechanism describes
the actual offer assignment, customer selection and outcome clock.
\citep{exec17_benchmark_cello}

Table 2 reports GPT-4o overall accuracy increasing from 0.59 to 0.70 with
CELLO-CoT, while confounder identification increases from 0.29 to 0.35.
Task-specific behavior matters: an aggregate improvement does not warrant
using a system to select adjustment variables. There is also an unresolved
arithmetic boundary in the published table. The twelve displayed
GPT-4o-plus-CELLO-CoT task entries are
\[
 0.83,\ 0.70,\ 0.52,\ 0.50,\ 0.35,\ 0.75,
 \quad 0.81,\ 0.72,\ 0.79,\ 0.77,\ 0.90,\ 0.61.
\]
Their equal-task mean is 0.6875. Under conventional rounding to two decimal
places, even a conservative closed envelope for that mean is
\begin{equation}
 [0.6825,\,0.6925]\ \cap\ [0.695,\,0.705]=\varnothing,
 \label{exec17:benchmark-rounding}
\end{equation}
where the second interval is the corresponding envelope for the reported
overall 0.70. With exactly 100 binary correctness outcomes per task, the
displayed entries would sum to 825 correct answers out of 1,200. This is a
conditional aggregation check against the published original. The cause of
the difference is unknown; no sample-level result or alternative aggregation
has been reconstructed. Consequently, the reported improvement remains an
author report with this unresolved qualification, rather than an independently
reproduced effect of the prompting procedure.
\citep{exec17_benchmark_cello,exec17_benchmark_exact}

\subsection{Refusal can be correct for the wrong reason}

CELLO's robustness exercise reformulates the questions as requests with
additional premises and binary answers. Section 6.4 states that the correct
answer is consistently ``No.'' This is a separate distribution from the
nearly balanced original binary questions. A constant-No classifier scores
100\% on the reformulated test without inspecting the scene. The paper also
describes GPT-4o declining requests because it is a language model unable to
interact with the physical world, rather than explaining the scene's causal
constraints. Correct answer polarity can therefore conceal a mismatch
between the tested mechanism and the reason for refusal.
\citep{exec17_benchmark_cello}

For a banking assistant, ``I cannot execute a payment'' and ``the available
evidence does not identify the effect of this payment policy'' have different
meanings. The first states an action limitation; the second states an
information limitation. A system could utter either sentence regardless of
whether a proposed comparison is scientifically valid. An evaluation should
therefore record the answer, the asserted reason and the consequence for the
requested decision. It should include supported requests as well as
unsupported requests, and distinguish the two kinds of refusal. Otherwise a
system that rejects every causal analysis may appear dependable while
providing no usable analysis.

Publicly constructed paired cases offer an inexpensive way to test this
distinction before commissioning a bank experiment. Begin with a mechanism
whose query is identified under the supplied assumptions. Create a companion
case that removes the assignment information or introduces an unresolved
alternative mechanism while keeping the business language similar. The
supported case should receive the warranted query or bounded answer; the
unsupported case should preserve uncertainty and name the missing information.
A complementary pair can retain the information and change only whether the
assistant has authority to act. This tests whether the assistant confuses
execution authority with causal identification. These are proposed evaluation
contrasts. Their merit is that the mechanism and the changed information are
explicit; a public observational label alone would not supply their causal
truth.

\subsection{Lexical agreement and directional reasoning}

CausalReasonBench supplies twelve underlying scenarios spanning four domains
and three difficulties. It repeats them over four task types, four prompting
strategies, three models and two replicates, yielding 1,152 responses, or 384
per model. The models are the paper's Groq-hosted LLaMA 3.1-8B, LLaMA 3.3-70B
and Qwen3-32B settings. The main settings use temperature zero, with the initial
contrast-case pass at 0.1 and seed 42 where supported. These details delimit a
historical evaluation; the response count is not a count of independently
sampled causal mechanisms. The gold explanations are author-defined using
the stated reference sources. In the minimum-wage example, the appendix itself
notes empirical debate. Following a stipulated textbook mechanism is a
different task from resolving that debate from an observational
vignette.~\citep{exec17_benchmark_cherif}

The paper's Causal Identification Rate, abbreviated CIR, is a token-set F1
score. For response-token set $R$, gold-keyword set $G$ and
$\varepsilon=10^{-8}$, its precision and recall are
$P=|R\cap G|/(|R|+\varepsilon)$ and
$Q=|R\cap G|/(|G|+\varepsilon)$, and
$\mathrm{CIR}=2PQ/(P+Q+\varepsilon)$. Causal Logic Precision, CLP, separately
assigns 1 to a gold forward pattern, 0.5 to ambiguous directional language,
and 0 to a reverse or absent pattern. Counterfactual Coherence Ratio, CCR,
uses counterfactual-token F1 for counterfactual tasks and equals CIR otherwise.
Thus the composite score outside the counterfactual tasks reduces to
\begin{equation}
 \mathrm{CAR}
 =0.40\mathrm{CIR}+0.35\mathrm{CLP}+0.25\mathrm{CCR}
 =0.65\mathrm{CIR}+0.35\mathrm{CLP}.
 \label{exec17:benchmark-score}
\end{equation}
The apparent three-component assessment has two distinct components in that
regime. This follows directly from the published definitions; it is not an
empirical claim about the models.
\citep{exec17_benchmark_cherif,exec17_benchmark_exact}

Set overlap ignores token order. For example, ``funding causes deposits'' and
``deposits causes funding'' give the same token set and hence the same CIR
against that gold set. The deliberately terse reverse statement isolates the
set operation; grammatical correction is a separate language task. CLP is
intended to test direction and could distinguish the two. The paper describes
whitespace tokenization, case normalization, Porter stemming and directional
regular expressions with up to five intervening words. Its complete executable
patterns have not been inspected here. The example therefore proves a limit
of the CIR component, not an exploit of the implementation or an equal score
under the whole composite. A banking evaluation should score the proposed
cause, outcome, direction, conditioning information and intervention
separately, while retaining the free-text explanation for review.
\citep{exec17_benchmark_cherif,exec17_benchmark_exact}

The paper selects the composite weights from fifteen candidates using 24
responses rated by two experts. The reported selected Kendall correlation
of 0.78 describes agreement on the selection set. It does not provide a
held-out validation of the selected metric. The required next evidence is a
fixed metric applied to different mechanisms and independently adjudicated
responses, including plausible keyword-rich errors, appropriate refusals and
alternative correct phrasings. This protects against a score that rewards
similar vocabulary while missing the banking decision's meaning.
\citep{exec17_benchmark_cherif}

Reproduction is also a substantive limit. The methods describe code, prompts,
gold annotations and raw results as released publicly, whereas the Data
Availability Statement on physical page 33 says they are available on
reasonable request. No public executable copy has been established by this
inspection. Both statements must be retained; the paper alone cannot qualify
a runnable integration. Its historical zero API charge under a free tier and
roughly eighteen hours of runtime likewise leave annotation, review, waiting,
infrastructure and current prices unmeasured. These limits support using the
paper as a source of test-design ideas, while leaving direct adoption and full
comparative economics unresolved.~\citep{exec17_benchmark_cherif}

\subsection{The monetary meaning of answering and withholding}

An all-No benchmark rewards withholding regardless of its opportunity cost.
To expose that omission, consider a deliberately stipulated banking review
task. A request is either supported or unsupported by its supplied information.
Endorsing a supported request earns USD 20; endorsing an unsupported request
loses USD 80; withholding earns zero. These amounts represent the value and
error cost of a review decision, not the deposit offer's actual causal effect
or an observed bank return. Let $\rho$ be the supported fraction, $t$ the
probability of endorsement given support and $f$ the probability of endorsement
without support. Then
\begin{align}
 A&=\rho t+(1-\rho)(1-f),
       \label{exec17:benchmark-accuracy}\\
 V&=20\rho t-80(1-\rho)f.
       \label{exec17:benchmark-value}
\end{align}
Here $A$ is classification accuracy and $V$ is expected USD per request.
They answer different questions. At the balanced mix $\rho=1/2$, the value
reduces to $10t-40f$.
\citep{exec17_benchmark_exact}

\begin{table}[htbp]
\centering
\caption{Exact illustrative support-review decisions at a balanced case mix.
All rates and payoffs are stipulated.}
\label{exec17:benchmark-bank-table}
\begin{tabular}{lrrrr}
\toprule
Rule & $t$ & $f$ & Accuracy & USD/request\\
\midrule
A & 0.80 & 0.10 & 0.850 & 4.00\\
B & 0.95 & 0.20 & 0.875 & 1.50\\
Always withhold & 0 & 0 & 0.500 & 0\\
\bottomrule
\end{tabular}
\end{table}

Rule B has higher accuracy and lower monetary value than rule A. A separate
400-request integer enumeration, with 200 supported and 200 unsupported
requests, gives A 160 true and 20 false endorsements, and B 190 true and 40
false endorsements. The resulting totals are USD 1,600 and USD 600,
respectively, reproducing the per-request values in the table. Always
withholding earns zero and obtains 50\% balanced accuracy, although it would
obtain 100\% accuracy if every label were unsupported. The comparison explains
why the refusal benchmark's label mix and the banking loss function must
remain visible.
\citep{exec17_benchmark_exact}

For a calibrated probability $p$ that a request is supported, endorsement has
conditional value $20p-80(1-p)$. It beats zero-value withholding only for
$p>0.8$, with indifference at 0.8. Calibration here concerns the support-review
label in its specified population. It does not calibrate an intervention
effect, a confidence interval or an assistant's self-reported confidence.
Changing the case mix or loss changes the decision problem. A public prototype
can vary these declared quantities to expose an unfavorable tradeoff; it
cannot discover the bank's actual losses merely by computing this threshold.
\citep{exec17_benchmark_exact}

\subsection{Repeated prompts are not independent mechanisms}

Prompting strategies should ordinarily be compared on the same mechanisms,
because differences in case difficulty otherwise obscure the contrast. Let
$D_{ij}$ be the paired utility difference between two strategies on mechanism
$i$ and within-mechanism variant $j$. An illustrative random-effects law is
\begin{equation}
 D_{ij}=\mu+U_i+E_{ij},\qquad
 \operatorname{Var}(U_i)=\tau^2,\quad
 \operatorname{Var}(E_{ij})=\sigma^2,
 \label{exec17:benchmark-cluster-law}
\end{equation}
with independent mechanisms, independent idiosyncratic terms and all such
terms mutually independent. Shared $U_i$ makes variants of the same mechanism
dependent even when their generation calls are separate. With $n$ mechanisms
and $r$ equally weighted variants per mechanism, the grand mean has variance
\begin{equation}
 \operatorname{Var}(\overline D)
 =\frac{\tau^2+\sigma^2/r}{n},\qquad
 \frac{\operatorname{Var}(\overline D)}{(\tau^2+\sigma^2)/(nr)}
 =1+(r-1)\frac{\tau^2}{\tau^2+\sigma^2}.
 \label{exec17:benchmark-cluster-variance}
\end{equation}
The denominator in the ratio is the variance that would follow from falsely
treating the variants as independent. For $n=12$, $r=8$ and
$\tau^2=\sigma^2=1$, the true variance is $3/32$, the independent-case
calculation gives $1/48$, and their ratio is $4.5$. Counting every ordered
within-mechanism covariance gives the same result. These are exact deductions
under the stated illustrative law, not estimates of dependence in either
paper. In particular, this calculation does not replace or reconstruct the
published significance tests or bootstrap intervals.
\citep{exec17_benchmark_exact}

Extra variants reduce the idiosyncratic contribution but do not remove the
mechanism contribution. A proposed prototype should therefore preserve the
mechanism identifier, compare strategies in pairs and allocate additional
budget to new mechanisms when mechanism variation dominates. With very few
mechanisms, simply choosing a cluster bootstrap does not guarantee reliable
uncertainty. The design must state its sampling population, weighting and
inferential assumptions, and may need to report a finite-case comparison
without a population claim. None of the twelve source scenarios becomes an
independent bank population because many prompts were generated from it.

These two benchmarks suggest complementary roles in the proposal. CELLO
offers a structured translation from entities through graphs to questions,
together with a warning that refusal polarity can conceal a semantic failure.
CausalReasonBench offers explicit automatic-score definitions and prompting
contrasts, together with limits in token matching, metric selection, repeated
mechanisms and reproducibility. The proposed evaluation should adopt those
test ideas only with explicit labels and information regimes, constant-answer
comparators, semantic adjudication, paired mechanisms and declared review
losses. Public cases and synthetic mechanisms can establish whether a narrow
prototype is worth further investigation. Authentic assignment information,
bank outcomes and institutional economics remain different evidence.
\citep{exec17_benchmark_cello,exec17_benchmark_cherif}
\label{exec17:benchmark-development-end}


% SOURCE docs/release/2026-10-29/source/docs/writing-repair/2026-10-29/src/c3364f0f26f32be5.tex CHAR 0
\section{What literature disagreement can establish about a benchmark}
\label{exec17:validity-development}

A benchmark graph is a scientific object with a population, measurement scheme
and intended regime. Its usefulness does not follow from how often it has been
reused. Yet a new paper that appears to contradict an arrow is not automatically
a correction to that graph. In banking, an association between a credit offer
and repayment can concern applicants, approved borrowers or customers who
accepted the offer. A different selection rule changes the population law.
An account-level measurement can also differ from a customer-level mechanism.
Before an assistant votes for or against a diagram, it must establish what
proposition the diagram and each retrieved result actually assert.

Zhang and colleagues offer a useful starting point: retrieve domain papers,
derive association queries from a benchmark graph, and ask a language model
whether each paper supplies relevant empirical evidence. Their June 2026
preprint reports a survey of 96 papers using 69 benchmarks, followed by an
evaluation of 11 graphs using 38,081 full texts. The approach makes benchmark
assumptions objects of investigation rather than treating every inherited
diagram as unquestionable truth. Its statistical and graphical interfaces
nevertheless need repair before the proposed library could use the result to
change a causal model. The assessment here concerns the inspected first-version
text, pseudocode and prompt; it does not claim to reproduce the upstream
implementation.~\citep{exec17_benchmark_consistency}

\subsection{A graph implication is not its converse}

For a distribution $P$ that is Markov with respect to a DAG $G$, graphical
separation supplies an independence implication:
\begin{equation}
 X\perp_G Y\mid Z \quad\Longrightarrow\quad X\perp_P Y\mid Z.
\label{exec17:validity-markov}
\end{equation}
Here $\perp_G$ denotes d-separation and $\perp_P$ statistical conditional
independence. Reversing this implication requires a further assumption, such
as faithfulness: the distribution has no extra independences beyond those
implied by the graph. Even dependence need not be captured by a particular
linear correlation coefficient. The paper's description that unseparated
variables, including adjacent variables, are always correlated therefore goes
beyond the Markov premise.~\citep{exec17_benchmark_consistency}

An exact project example shows that positive density does not repair the
converse. Let independent disturbances $\varepsilon_X,\varepsilon_A,
\varepsilon_Y$ be standard Gaussian and define
\begin{equation}
 X=\varepsilon_X,\qquad A=X+\varepsilon_A,\qquad
 Y=X-A+\varepsilon_Y.
\label{exec17:validity-cancellation}
\end{equation}
The graph contains $X\to A$, $X\to Y$ and $A\to Y$, and the joint law is
Markov for this graph. In the order $(X,A,Y)$ its covariance matrix is
\begin{equation}
 \Sigma=\begin{pmatrix}1&1&0\\1&2&-1\\0&-1&2\end{pmatrix},
 \qquad \det(\Sigma)=1,\qquad
 \operatorname{Cov}(X,Y\mid A)=\tfrac12.
\label{exec17:validity-covariance}
\end{equation}
Indeed $Y=-\varepsilon_A+\varepsilon_Y$, so $X$ and $Y$ are marginally
independent despite their adjacency. The independent-disturbance coefficient
map is nonsingular, establishing positive density. The direct and mediated
effects cancel for this population; conditional association need not cancel.
The covariance and determinant are reproduced by exact coefficient arithmetic,
not by a sample whose correlations happen to be small.
\citep{exec17_advice_validity_exact}

For an offer, subsequent account action and final outcome, such a construction
distinguishes a mechanism from its total marginal association. It is not an
assertion that the bank has these coefficients. Its role is to show why a
literature result reporting independence cannot automatically invalidate an
adjacency under a Markov-only contract. A faithful-model comparison is a
different, explicitly stronger exercise. Conversely, an estimated nonzero
correlation does not by itself identify a direct arrow. A benchmark evaluator
must label which of these propositions it tests.

\subsection{Conditioning can open a path that a hitting set overlooks}

The published Algorithm 1 enumerates collider-free paths and returns minimal
sets that intersect each such path as conditional-independence sets. Blocking
paths that are initially active is necessary, but does not ensure that
conditioning leaves every other path blocked. A collider becomes active when
it, or a descendant, is conditioned upon. The paper's own d-separation
definition includes this condition; the inspected hitting-set pseudocode does
not apply it to candidate separators.~\citep{exec17_benchmark_consistency}

Consider the following five-node project DAG:
\begin{equation}
 X\to C,\quad B\to C,\quad C\to Z,\quad X\to Z,
 \quad Z\to Y,\quad B\to Y.
\label{exec17:validity-graph}
\end{equation}
Its collider-free paths from $X$ to $Y$ are $X-Z-Y$ and $X-C-Z-Y$. The unique
minimal hitting set of these paths is $\{Z\}$. Conditioning on $Z$ blocks
both, but activates $X\to C\leftarrow B\to Y$ through the descendant
$Z$ of collider $C$. It also activates the path
$X\to Z\leftarrow C\leftarrow B\to Y$ by conditioning on collider $Z$
itself. Thus $\{Z\}$ is not a d-separator. Enumerating every simple path
and applying the full activity criterion gives $\{B,Z\}$ as the minimal
d-separator. This is a counterexample to the published specification, not an
observed output of uninspected code.~\citep{exec17_advice_validity_exact}

A nonsingular Gaussian construction verifies the practical consequence:
\begin{equation}
\begin{split}
 X&=\varepsilon_X,\qquad B=\varepsilon_B,\qquad
 C=X+B+\varepsilon_C,\\
 Z&=X+C+\varepsilon_Z,\qquad Y=Z+B+\varepsilon_Y,
\end{split}
\label{exec17:validity-opened-scm}
\end{equation}
with independent standard Gaussian disturbances. Its coefficients give
$\operatorname{Var}(Z)=7$, $\operatorname{Cov}(X,Y)=2$,
$\operatorname{Cov}(X,Z)=2$ and $\operatorname{Cov}(Z,Y)=8$. Therefore
\begin{equation}
 \operatorname{Cov}(X,Y\mid Z)=2-\frac{2\times8}{7}=-\frac27.
\label{exec17:validity-opened-covariance}
\end{equation}
The exact covariance agrees with the independent active-path enumeration.
\citep{exec17_advice_validity_exact}

The remedy is concrete. Candidate conditioning sets must pass a sound
d-separation procedure for the stated graph family, including collider and
descendant activation. Small complete path enumeration is an independent
reference for testing a scalable implementation; it is not a scalable method
for every bank graph. If a graph has latent variables or selection nodes, the
evaluator needs the appropriate representation and separation criterion,
rather than treating its observed mixed graph as an ordinary DAG. The query
also needs to distinguish conditioning on a recorded covariate from
intervening on it. A controlled experiment and an observational subgroup do
not generally describe the same probability distribution.

\subsection{A nonsignificant test does not certify independence}

The inspected prompt instructs the verifier to infer independence when a
correlation test is not significant. That rule confuses a sample decision
with a population property. It can turn an underpowered study into apparent
evidence against a correct causal mechanism. It also treats a selected test
as if it were sensitive to every form of dependence.
\citep{exec17_benchmark_consistency}

For a finite illustration, let $X$ be a randomized binary treatment with
probability $1/2$ and let an independent flip $V$ have probability $1/4$.
Set $Y=X\mathbin{\mathrm{xor}}V$. The four population cell probabilities are
\begin{equation}
 (p_{00},p_{01},p_{10},p_{11})=\tfrac18(3,1,1,3),
 \qquad E{Y\mid\operatorname{do}(X=1)}
       -E{Y\mid\operatorname{do}(X=0)}=\tfrac12.
\label{exec17:validity-binary}
\end{equation}
In four independent observations, the event of one observation in every cell
has probability
\begin{equation}
 4!\left(\frac38\right)^2\left(\frac18\right)^2
       =\frac{27}{512}.
\label{exec17:validity-balanced}
\end{equation}
For that balanced table, Fisher's two-sided exact test has $p=1$. Conditional
on both margins being $(2,2)$, the upper-left cell under independence takes
values $0,1,2$ with probabilities $(1,4,1)/6$; all three tables have probability
no greater than that of the observed central table. The resulting sum is one.
Complete enumeration of all 256 four-observation sequences finds the 24
balanced sequences and the probability above. Yet the true randomized effect
remains $1/2$.~\citep{exec17_advice_validity_exact}

If an outcome contributes a stipulated USD 100 and treatment costs USD 30, the
population incremental value is $100(1/2)-30=\text{USD }20$ per decision.
Neither the four-row sample nor its large $p$ value establishes this value; it
is known here only because the example stipulates the structural law. The
example demonstrates the possible loss from converting a nonsignificant
finding into a zero-effect declaration. It is not a proposal ROI estimate.

A repaired extraction should preserve the estimator, sample size, sampling
design, uncertainty interval, null hypothesis and covariates. ``Did not reject
this null'' is an admissible extracted finding. ``Established independence''
requires a different justification. For a decision, an equivalence analysis
could target a scientifically chosen negligible-effect region, with its own
design and error control. A language model may help locate those details; it
cannot supply missing power, overlap or identification by rewriting the prose.

\subsection{Unknown answers change the target of an inconsistency rate}

The paper retains an unknown category and excludes unknown answers from its
inconsistency denominator. This is a useful distinction from forcing every
paper into a binary label, but it changes what a reported rate means. Write
$R=1$ when a query receives an answer and $D=1$ when that answer is categorized
as a contradiction. If $r=P(R=1)$ and $q=P(D=1\mid R=1)$, then without an
assumption on unknown cases, a latent full-query contradiction fraction $d$
obeys only
\begin{equation}
 rq\leq d\leq rq+1-r.
\label{exec17:validity-unknown}
\end{equation}
The bound assumes that a meaningful binary contradiction status exists for
every target query and that answered labels are correct; context mismatch or
misclassification requires further treatment. The original answered-case
metric does not make either additional premise true.
\citep{exec17_benchmark_consistency}

In an authored example with 20 contradictions among 30 answered queries out
of 100, the answered-case rate is $2/3$, the answer rate is $0.3$, and the
full-query range is $[0.2,0.9]$. Both endpoints are attained by assigning all
unknown cases to one category. This is a missing-information bound, not an
interval for statistical sampling error.~\citep{exec17_advice_validity_exact}
The pipeline should therefore report answer coverage alongside disagreement,
and distinguish unique papers, paper-pair queries and graph propositions.
Several queries can share one paper, and several papers can share a dataset;
their labels are not automatically independent replications.

Retrieval introduces another selection layer. The inspected procedure uses
PubMed/PMC and OpenAlex-related access routes, caps ranked search results at
500 per pair, and caps examined pairs at 200. Its verifier is the historical
Qwen3-30B-A3B-Thinking-2507 model, reported on an H100 80GB. Accessible text,
search ranking, variable synonyms and the relevance filter define the resulting
evidence pool. An inconsistency rate in that pool is not a random-sample
estimate of all scientific evidence or all graph errors. The authors report
90 correct judgments in a manual sample of 100; the linked detailed artifact
has not been inspected here, so this remains their reported validation.
\citep{exec17_benchmark_consistency}

There are also unresolved source-accounting questions. The main text lists
69 benchmarks while the appendix enumeration reaches 72. The introduction
describes Barley as misaligned, but the evaluated set excludes it. A separate
bibliometric search reports 11,633 LLM causal-discovery papers in 2025, while
its appendix includes MATH, ImageNet and COCO among prominent benchmark
labels. Without resolving search specificity, that count cannot establish
the size of the causal-discovery field. These limitations do not erase the
value of evidence retrieval; they limit the numerical and substantive claims
that this proposal can inherit.~\citep{exec17_benchmark_consistency}

\subsection{A repair that retains the useful human workflow}

The useful adoption target is a reviewed evidence queue. Each item should link
an exact passage to the original PDF and record population, unit, variable
definition, measurement date, conditioning set, intervention and selection
mechanism. A proposed conflict should identify both propositions and explain
why they refer to a comparable regime. For a banking diagram, a relationship
among approved borrowers cannot silently replace one among all applicants.
An updated sensor or accounting definition can require a new variable and
measurement model rather than an arrow reversal.

The proposal can conditionally adopt retrieval, synonym suggestions and
abstaining evidence extraction. It should reject automatic benchmark relabeling
from literature votes, adjacency-as-correlation rules and nonsignificance-as-
independence prompts. Direct software adoption or a port remains contingent on
inspecting pinned code, dependencies and rights. The mathematical repairs are
known: full separation checks, faithful versus Markov-only query types, correct
statistical statements, explicit missing-answer targets and reviewed context
alignment. Whether a repaired pipeline attains useful accuracy and cost is a
prototype question, not a result established by the present manuscript.

A public prototype can make that question testable before a bank trial. Build
a small adjudicated set of original passages and graph queries, including
known collider cases, cancellation, nonsignificant dependent samples, changed
populations and deliberately unavailable evidence. Keep both the graph truth
for constructed cases and the expert justification for real passages separate
from the model's suggestions. Compare keyword retrieval, human-only review and
an abstaining assistant on factual extraction, query validity, answer coverage,
false conflict alerts, review time and decision-relevant consequences. Preserve
disagreements instead of training on the test judgments. The mathematical
examples specify failures that can be tested now; broader scientific truth
and banking relevance require the substantive evidence that the workflow is
designed to help a reviewer inspect.

\label{exec17:validity-development-end}


% SOURCE docs/release/2026-10-29/source/docs/writing-repair/2026-10-29/src/24b3b0d04fd78fd6.tex CHAR 0
\section{Evaluate a causal coding workflow beyond its test-pass rate}
\label{exec19:evalcode-scope}

A causal-analysis assistant can generate syntactically correct code, execute a
query, and pass tests while answering the wrong banking question. A repayment
model might use postdecision information; an effect estimator might target the
wrong cohort; a reporting agent might convert an execution failure to a zero
effect. Evaluation must therefore identify the task, the observations on which
performance is measured, and the protocol that decides whether an output is
usable. Those choices connect the software interface to the scientific target
and to the economic case for a prototype.

Two selected historical surveys help organize this distinction. Chang and
coauthors' evaluation survey, inspected at its 6 July 2023 initial version,
separates what to evaluate, where to evaluate and how to evaluate. Its summary
explicitly limits its success and failure conclusions to particular datasets.
It discusses automated and human assessment, prompt robustness, complete
behavioral tests and evolving benchmarks. It also asks whether a dynamically
constructed test is truly out of distribution. Dynamic generation is thus a
design option, not a certificate of novelty or freedom from contamination.
These are historical organizing arguments, not a current ranking of models.
\citep{exec19_surveys_evaluation}

Jiang and coauthors' code-generation survey, inspected at its 1 June 2024
initial version, covers data curation, model adaptation, repository and agent
workflows, evaluation and applications. It distinguishes text/structural
similarity, execution metrics, human assessment and model judging. The inspected
performance tables compile cited model results; they are not a harmonized
experiment reproduced here. Its framework descriptions likewise need the
constituent technical originals before a package or claimed score can enter
our admitted design. The selected PDF's template footer says September 2024,
whereas the archive and first-page stamp identify the June initial version.
That footer does not establish journal publication or equivalence to the later
record that led to discovery.~\citep{exec19_surveys_code}

\subsection{A metric must define its sampling pool and selector}
\label{exec19:evalcode-selection}

The code survey's equation for pass-at-$k$ uses a finite-pool expression of the
form $1-\binom{n-c}{k}/\binom{n}{k}$. Its surrounding prose, however, describes
$c$ as the number of correct samples within the selected $k$, despite defining
$n$ as the larger generated pool. These definitions conflict. Before adopting
the expression, the implementation must fix which pool the count refers to.
The following derivation is our own finite counting argument, not an attribution
of an uninspected original theorem.~\citep{exec19_surveys_code}

Let a fixed pool contain $n=10$ candidates, of which exactly $c=3$ pass a
stipulated test suite. Choose $k=2$ candidates uniformly without replacement.
There are $\binom{10}{2}=45$ pairs, of which $\binom72=21$ contain no passing
candidate. Thus
\begin{equation}
 P(\hbox{at least one passes})
 =1-\frac{\binom72}{\binom{10}2}=\frac{24}{45}=\frac8{15}.
 \label{exec19:evalcode-pool}
\end{equation}
The successful count in that formula is the number among all ten candidates.
Substituting the random number of successes in the selected pair instead gives
a different quantity; exact enumeration gives average $16/135$, not $8/15$.
The source inconsistency must therefore be resolved before it becomes a
benchmark implementation or an accepted numerical comparison.
\citep{exec19_surveys_reference}

An at-least-one success event does not specify how the deployed system finds
that success. If it chooses one of the two sampled candidates uniformly,
its chance of selecting a passing candidate remains $3/10$. Reaching $8/15$
as an output-success rate requires an oracle or a suitably qualified selector
that can recognize the passing candidate whenever one is present. A selector
using public tests, a model judge or analyst review has its own errors and
costs. For causal code, even a perfect test-suite selector establishes only
success under those tests; it does not establish identification assumptions or
the relevance of the estimand.~\citep{exec19_surveys_reference}

\subsection{Retain the attempted-request denominator}
\label{exec19:evalcode-pipeline}

Consider a second authored example with twenty attempted causal coding tasks.
Sixteen produce runnable outputs; twelve pass their unit tests; seven preserve
the declared causal target and information clock; five satisfy the full output
contract, including evidence and uncertainty. These nested counts are stipulated
to illustrate measurement design, not observations from a model. The two
headline rates are different:
\begin{equation}
 r_{\rm tests\mid runnable}=\frac{12}{16}=\frac34,
 \qquad
 r_{\rm accepted\mid attempted}=\frac5{20}=\frac14.
 \label{exec19:evalcode-denominators}
\end{equation}
A report containing only the first rate would omit failed execution and
scientific-contract failures. The pipeline should retain the request identifier,
eligibility decision, every candidate and retry, selected result, failure or
abstention reason, and final adjudication. The rates need not be multiplied
under an independence assumption: the example directly counts nested outcomes.
Correlated prompts and shared graph families also mean these twenty rows do
not automatically justify a binomial interval or twenty independent units of
evidence.~\citep{exec19_surveys_reference}

This requirement follows the surveys' limitations rather than treating a
benchmark as a universal standard. The code survey stresses that execution
metrics depend on unit-test quality and omit attributes such as maintainability,
efficiency, stability and repository dependencies. Human review adds expertise
but also cost, variance and reproducibility issues. Model judges can inherit
position, verbosity, self-enhancement and reasoning biases. The evaluation
survey similarly treats evaluator variability and the validity of the testing
system as questions requiring scrutiny. A causal workflow adds its own
target, graph, clock, support and refusal semantics to this broader software
assessment.~\citep{exec19_surveys_code,exec19_surveys_evaluation}

\subsection{Count costs on all attempts and benefits on accepted outputs}
\label{exec19:evalcode-cost}

Suppose generation costs USD 2 and the test harness costs USD 1 for every
attempt, while specialist review costs USD 5 for each test-passing output.
Suppose an accepted output is worth USD 30 under an explicitly hypothetical
valuation. For the twenty-task construction,
\begin{align}
 C&=20(2+1)+12(5)=120\quad\text{USD},
 \label{exec19:evalcode-costs}\\
 C/N_{\rm accepted}&=120/5=24\quad\text{USD per accepted output},
 \label{exec19:evalcode-unitcost}\\
 B-C&=5(30)-120=30\quad\text{USD in total}.
 \label{exec19:evalcode-net}
\end{align}
Net value is USD 1.50 per attempted task. Equating all twelve test passes with
accepted outputs would instead produce USD 240 net and overstate value by
USD 210. These prices and benefits are assumptions; none is a model-provider
quote, observed analyst saving or bank ROI. The example demonstrates why the
all-attempt cost ledger and acceptance rubric must be frozen together before
a prototype trial. A real experiment should additionally include setup, retries,
timeouts, correction, dependency maintenance and the cost of rejected work.
\citep{exec19_surveys_reference}

\subsection{A public evaluation design for the proposed causal assistant}

The proposal can test these ideas without authentic bank records. Generate
declared structural and tabular examples with known interventions, ambiguous
graphs, missing outcomes, invalid clocks and unsupported requests. Separate
development graph families from evaluation families and retain original
generation seeds and source versions. Score formulation, diagram semantics,
identification, code execution, numerical output, uncertainty, reporting and
appropriate refusal separately. The classifier or judge for each stage must
itself be checked on counterexamples; agreement with a fluent explanation is
not an independent causal oracle.

Use unit tests for declared numerical and interface properties, source-grounded
checks for algorithm assumptions, and adjudicated scientific review where the
target is ambiguous. Report automatic, human and model-judge decisions with
their disagreement rather than substituting one for another invisibly.
Randomize order when testing a model judge, retain its version and prompt, and
evaluate selection separately from candidate generation. The survey discussions
motivate these dimensions but do not prove that this proposed protocol is
calibrated. Historical benchmarks can supply leads and adverse examples;
current package adoption requires its own pinned originals and conformance.
\citep{exec19_surveys_evaluation,exec19_surveys_code}

The next experiment is consequently a bounded test of whether the assistant
preserves a declared causal contract at a useful all-request cost. Public and
synthetic tasks can expose wrong denominators, selection optimism, unsafe
conversions and semantic errors before costly bank experiments. They cannot
verify an institution's actual assignment mechanism, audit truth or realized
economic return. Those remain later questions after the prototype demonstrates
why further evidence would be worth acquiring.


% SOURCE docs/release/2026-10-29/source/docs/writing-repair/2026-10-29/src/854592282bfc9971.tex CHAR 0
\section{Retrieval and adaptive repair need separate evidence and acceptance contracts}
\label{exec19:ragrepair-scope}
A causal banking assistant needs to retrieve the right evidence, translate it
into a defensible intervention question, produce an executable analysis, and
recognize when the result does not answer that question. Retrieval and iterative
repair can support these stages, but their success criteria differ. A faithfully
quoted passage can be scientifically inapplicable; a program can pass the tests
it has repeatedly seen while using the wrong time window or treatment. The
proposal must therefore preserve the chain from source identity through causal
meaning to held-out decision quality.

Two historical originals help sharpen this design. Gao and colleagues' survey
is inspected at arXiv version 5 of 27 March 2024, distinct from its December 2023
initial discovery date. The CoT-SelfEvolve study is inspected at arXiv version 1
of 28 August 2024. Its first page and running footer also contain a 2025 conference
template and placeholder DOI; those do not establish a verified proceedings
version. Archive metadata lists its first author as Quoc, Thai Tang, while the
PDF displays Thai Tang Quoc. The version and author-name distinction is retained
in the audit. Neither original establishes complete October 2026 frontier coverage.
\citep{exec19_rag_survey,exec19_selfcorrection_study}

\subsection{Retrieval faithfulness is narrower than causal evidence}
\label{exec19:ragrepair-retrieval}
The survey distinguishes a basic index--retrieve--generate pipeline, advanced
pre- and post-retrieval improvements, and modular pipelines that can rearrange
or repeat retrieval and generation. Basic indexing converts source documents
into chunks and representations; retrieval ranks context for a query; generation
uses that context to answer. Query rewriting, finer indexing, metadata, reranking
and compression can improve the fit between question and material. Modular
search and iterative retrieval provide alternatives when a single initial
query is inadequate. These are useful architectural choices, not proof that a
retrieved statement is true or that a generated graph is identified.
\citep{exec19_rag_survey}

For a bank asking whether a servicing call prevents default, retrieval must
retain more than a passage containing both terms. An observational prediction
study, a randomized call policy and a vendor's marketing description can all
be semantically similar to the question while supporting different claims.
The proposed evidence object therefore records source type, original date and
version, population, treatment and outcome definitions, assignment mechanism,
measurement clock, assumptions and the precise claim supported. Rewriting the
business question may produce several candidate estimands; it must not silently
replace an intervention with a correlation because that version retrieves more
text. This evidence-object design is a project proposal motivated by the survey's
retrieval and augmentation distinctions, not a capability reported for an
already integrated package.
\citep{exec19_rag_survey}

The survey separately evaluates retrieval quality and generation quality.
It distinguishes context relevance, answer faithfulness and answer relevance,
and discusses noise resistance, refusal when necessary knowledge is missing,
integration across documents and resistance to known false retrieved statements.
Its term counterfactual robustness means recognizing and rejecting inaccuracies
in documents. It does not mean identifying a potential outcome or preserving
cross-world structural dependence. Confusing these uses of the same word would
let a language benchmark appear to qualify a causal capability it never tested.
The survey also warns that the listed evaluation metrics were not yet a mature,
standardized system. A package name or metric family in its comparison is a
lead for primary inspection, not an automatically adopted scoring implementation.
\citep{exec19_rag_survey}

A public evaluation can make these distinctions concrete before a costly bank
trial. Construct matched evidence packets containing a relevant original,
a wrong-population study, a later revision, a contradictory source and an
unsupported confident claim. Preserve their identities and the gold-supported
conclusions. Score retrieval of the needed evidence separately from faithful
quotation, applicability to the declared estimand and justified refusal.
Then require the proposed graph to mark unsupported edges and alternatives.
Removing an irrelevant source should not alter a justified answer; removing
the only identification assumption should change its status. These are proposed
perturbation and evidence tests, not a claim that a retrieved paper establishes
an actual bank's assignment mechanism. The survey's discussion of noise,
contradiction and source/metadata exposure motivates such explicit checks.
\citep{exec19_rag_survey}

\subsection{What the self-correction study actually measures}
\label{exec19:ragrepair-study}
CoT-SelfEvolve combines retrieved StackOverflow discussions, generated reasoning
prompts, initial code, syntax checking, execution against unit tests and feedback-
guided revisions. Its feedback is partly external: tracebacks and test outcomes
come from an executor rather than from a model merely reconsidering its prose.
The inspected protocol uses the Completion setting of DS-1000, comprising
1,000 problems across seven Python libraries. It reports a retrieval collection
of 558,402 posts and 972,513 comments, with a minimum-comment filter and a token
budget for assembled documents. The paper also describes its benchmark tasks
as extracted from StackOverflow. This shared origin motivates an explicit
retrieval/benchmark overlap audit; the inspected descriptions do not establish
that leakage occurred or that every relevant overlap was removed.
\citep{exec19_selfcorrection_study}

The study defines pass@$n$ as the fraction of problems for which all tests pass
within at most $n$ sequential attempts, examining budgets from one to five.
Later attempts can use feedback from earlier ones. This differs from the event
that a uniformly chosen subset of an independent candidate pool contains a
passing program. It also differs from passing a fresh causal acceptance test.
The proposed product should retain the paper's operational definition rather
than treating identical-looking pass notation as a common metric.
\citep{exec19_selfcorrection_study}

The reported comparisons are informative but protocol-specific. Table 1 gives
an aggregate change from 30.5 to 35.5 percent when both prompt generators are
enabled; Table 2 gives 35.5 to 39.5 percent when the prompt generator uses the
larger model while the code generator remains fixed. Arithmetic on those
printed, rounded values gives
\begin{align}
 35.5-30.5&=5.0\text{ percentage points},&
 \frac{35.5-30.5}{30.5}&=\frac{10}{61},\nonumber\\
 39.5-35.5&=4.0\text{ percentage points},&
 \frac{39.5-35.5}{35.5}&=\frac{8}{71}.
 \label{exec19:ragrepair-percentages}
\end{align}
The relative changes are approximately 16.39 and 11.27 percent, respectively;
they are not increases of that many percentage points. Table 1 also shows
heterogeneity: the SciPy entry changes from 33.02 to 32.08 percent with both
prompt generators enabled. A positive aggregate therefore does not establish
a gain for every library or every causal-analysis task. No raw counts or paired
uncertainty interval is recovered from these rounded percentages.
\citep{exec19_selfcorrection_study,exec19_retrieval_repair_reference}

For its GPT-4 attempt-budget experiment, Table 3 reports overall success rising
from 14.0 to 83.2 percent between one and five allowed attempts. The text reports
prompt tokens rising from 2.54 million to 9.29 million, while completion tokens
do not rise at the same rate. This is evidence about the reported historical
protocol, not a present model benchmark or a full cost-benefit analysis.
Cross-figure and cross-table totals should not be pooled without reconciling
settings and runs. The authors themselves discuss proprietary model drift,
temperature and sampling settings, and the dependence on available unit tests.
Their analogy between five refinement steps and five usability testers does
not supply an optimal-stopping theorem for code or causal analysis.
\citep{exec19_selfcorrection_study}

Executable reuse has an additional unresolved prerequisite. The selected
original's reproducibility statement gives \texttt{ANONYMIZED-LINK} for its code
and prepared data. That is not a resolved repository identity, version or
license. It is therefore appropriate to retain the method as a research
comparator while deferring code adoption. A similarly named repository cannot
be substituted without establishing authorship and correspondence. Proposed
CI integration in the discussion also does not demonstrate that such an
integration has already been built here.
\citep{exec19_selfcorrection_study}

\subsection{Count failed attempts and accepted causal outputs}
\label{exec19:ragrepair-accounting}
A small authored example shows why the stopping and acceptance records matter.
Consider ten requests with a maximum of three attempts. Their first test-passing
attempts are
\begin{equation}
 (1,1,2,2,2,3,\varnothing,\varnothing,\varnothing,\varnothing),
 \label{exec19:ragrepair-attempts}
\end{equation}
where $\varnothing$ means that no candidate passes within the budget. The
cumulative counts of passing requests are two, five and six. The six successful
requests use eleven attempts in total; the four failures consume twelve more.
Thus all-request cost uses 23 attempts, not eleven. The average attempt count
is 2.3 across all requests, compared with $11/6$ if failures are omitted.
Suppose a separate causal review accepts only four of the six test-passing
outputs because two violate the frozen estimand or information clock. Test
success is $6/10$, while accepted causal output is $4/10$.
\citep{exec19_retrieval_repair_reference}

Assume USD 2 per attempt and USD 15 value per accepted output. The complete
cost and net value are
\begin{align}
 C&=2(23)=\mathrm{USD}\ 46,\nonumber\\
 C/4&=\mathrm{USD}\ 11.50\text{ per accepted output},\nonumber\\
 V_{\mathrm{accepted}}&=15(4)-46=\mathrm{USD}\ 14,\nonumber\\
 V_{\mathrm{test\ only}}&=15(6)-46=\mathrm{USD}\ 44.
 \label{exec19:ragrepair-cost}
\end{align}
Equating test passes with accepted causal outputs overstates this stipulated
net value by USD 30. The example does not measure actual token prices,
engineering labor, model performance or bank ROI. Its role is to enforce a
complete denominator before real economic evaluation, including unsuccessful
retrieval, repair, runtime and human adjudication.
\citep{exec19_retrieval_repair_reference}

\subsection{A public prototype protocol before paid trials}
\label{exec19:ragrepair-protocol}
The proposal can now specify a bounded public experiment. Freeze the source
corpus and exclusions, model/backend version, prompt and package environment,
maximum attempts and stopping rule. Compare one-shot generation, executor-
feedback repair and retrieval-assisted repair under the same request budget.
Keep the feedback tests distinct from an independent evaluation set that checks causal
target, time alignment, graph and identification scope, uncertainty and refusal.
Include deliberate missing-evidence cases, misleading retrieval, future leakage,
wrong treatment definitions and syntactically correct but causally invalid
programs. Log every retrieved source and every attempt; assess accepted outputs
and complete cost jointly. This design combines the survey's component-level
criteria with the empirical study's adaptive feedback structure without
borrowing either source's reported score as a product guarantee.
\citep{exec19_rag_survey,exec19_selfcorrection_study,exec19_retrieval_repair_reference}

Proceeding requires a genuine artifact for an executable CoT-SelfEvolve
comparison, inspected originals for any adopted retrieval/evaluation package,
and an independently specified reference answer for evaluation. Where these are absent,
use explicitly authored reference components and label them accordingly.
Public evidence can justify whether a prototype merits an expensive banking
experiment; it cannot establish institutional data truth or replace subsequent
owner acceptance. The historical findings narrow the design and expose concrete
failure modes while leaving broader frontier coverage and adoption open.
\citep{exec19_rag_survey,exec19_selfcorrection_study}


% SOURCE docs/release/2026-10-29/source/docs/writing-repair/2026-10-29/src/70273c19b4bfad55.tex CHAR 0
\section{COAT: eliciting measured factors without assuming causal truth}
\label{exec19:coat-scope}

Consider a servicing team with customer narratives, offer decisions and later
repayment outcomes. Its first difficulty may be defining usable variables:
does a phrase describe a pre-existing cash-flow shock, a customer's reaction to
the offer, or a retrospective explanation written after default? A language
model can suggest a factor and extract its value, but the suggestion, the
measurement and the causal role require different evidence. The selected
13 October 2025 revision of Liu and colleagues' COAT framework addresses this
interface between unstructured observations and causal analysis. It proposes
and parses factors, then uses conventional causal-discovery feedback to refine
them. The paper separates COAT-MB for Markov blankets, COAT-PAG for refining
partial ancestral graphs, and COAT-ADJ for adjustment questions. This is a
useful architecture for a banking prototype whose variables are not already
settled.~\citep{exec19_coat}

\subsection{What the representation assistant contributes}
\label{exec19:coat-representation}
In the paper's notation, $X$ is unstructured input, $\mathcal U$ is a set of
structured target variables and $W=h(X)$ is a proposed structured
representation. A factor instruction specifies a deterministic map from the
input to a numerical value. The initial sampling setup uses independently
drawn paired observations. COAT-MB proposes factors in rounds, parses their
values and updates a target's Markov blanket. Its feedback selects a cluster
where the existing factors leave high conditional entropy of the target. The
algorithm stops at a maximum number of rounds or an unchanged blanket.
These choices describe an adaptive search procedure; the stopping condition
does not itself certify that all relevant factors have been found.
\citep{exec19_coat}

For the bank, a factor definition therefore needs an observable interpretation,
an extraction rule, a timestamp and a missingness convention. ``Temporary
income interruption'' should not acquire a different meaning depending on the
desired intervention answer. If it is read from a note written after repayment,
its predictive usefulness does not make it a pre-offer confounder. Proposed
factors can be retained as hypotheses while their provenance and clock are
examined. An ordinary prediction improvement is useful feedback, but it does
not decide which variables are eligible adjustment controls.

The following project derivation makes an information limit explicit.
For a fixed mapping $W=h(X)$, if a
relevant variable $U$ is independent of $X$, it remains independent of $W$:
for any event $B$ of the representation,
\begin{equation}
 P(U=u,W\in B)
 =\sum_{x:h(x)\in B}P(U=u)P(X=x)
 =P(U=u)P(W\in B)
 \label{exec19:coat-information}
\end{equation}
in the discrete case. Naming a plausible hidden cause cannot create its
customer-specific measurement. Additional records or informative measurement
assumptions may help; the example is not a prohibition on extracting genuine
information already present in text. The paper's ``hidden world'' concerns
eliciting high-level factors from observations, not an unconditional ability
to identify every unobserved cause.~\citep{exec19_coat}

\subsection{The parsing assumption is scientifically substantive}
\label{exec19:coat-parsing}
The selected paper explicitly assumes faithful and Markov correspondence
between a graph and the factor distribution, with conditional independences
preserved after parsing. It acknowledges that annotation error can violate
this premise. High marginal accuracy is consequently insufficient to establish
that a parsed variable can replace its true counterpart in adjustment or
conditional-independence tests.~\citep{exec19_coat}

An exact project-authored mechanism illustrates the issue. Let $U$ be a fair
binary pre-offer cause, let the offer $A$ and repayment $Y$ be conditionally
independent given $U$, and specify
\begin{equation}
 P(A=1\mid U)=P(Y=1\mid U)=\tfrac14+\tfrac12U,
 \qquad W=U\mathbin{\oplus}E,
 \qquad E\sim\operatorname{Bernoulli}(1/4),
 \label{exec19:coat-noise-law}
\end{equation}
where $E$ is independent of the other structural randomizations and $\oplus$
denotes exclusive-or. There is no $A\to Y$ arrow: the true offer effect is
zero, and both interventional repayment risks equal $1/2$. The proxy $W$ is
75\% accurate. Nevertheless, conditional on either value of $W$, $U$ retains
variance $3/16$. The conditional covariance is
\begin{equation}
 \operatorname{Cov}(A,Y\mid W=w)
 =\tfrac14\operatorname{Var}(U\mid W=w)=\tfrac3{64}>0.
 \label{exec19:coat-covariance}
\end{equation}
The factor error has destroyed the conditional independence that held given
$U$. This is a deliberate violation of the paper's parsing assumption, not
evidence of failure when that assumption holds.~\citep{exec19_coat_reference}

The exact conditional repayment risks are $3/10,1/2,1/2,7/10$ in the order
$(a,w)=(0,0),(1,0),(0,1),(1,1)$. Standardizing over the fair proxy gives
\begin{equation}
 \sum_w P(Y=1\mid A=0,w)P(w)=\tfrac25,
 \qquad
 \sum_w P(Y=1\mid A=1,w)P(w)=\tfrac35.
 \label{exec19:coat-proxy-effect}
\end{equation}
The spurious effect is $1/5$. With an assumed USD 100 payoff per incremental
repayment and USD 10 offer cost, that calculation yields USD 10, whereas the
true zero-effect offer has net value USD $-10$. The positive joint law is
enumerated with rational arithmetic; no COAT run, empirical annotation rate
or actual bank benefit is being reported. The proposed implication is to
evaluate extraction errors by variable, outcome, treatment and time, rather
than accepting a single overall accuracy number as a causal qualification.
\citep{exec19_coat_reference}

\subsection{Separate iterative discovery from a finite stopping guarantee}
\label{exec19:coat-attempts}
COAT's ability definition posits a positive probability of proposing a useful
factor and a positive multiplicative reduction in residual information. Its
convergence discussion then uses these quantities to characterize successful
discovery rounds. The proof of Proposition 5 treats the number of useful rounds
as binomial and explicitly approximates that distribution by a Gaussian.
That calculation is informative about the proposed regime, but an
uncontrolled approximation cannot be adopted as an exact finite probability
certificate for a banking run. The model for repeated useful proposals also
needs to match an adaptive agent's history and failure modes.~\citep{exec19_coat}

For example, four independent proposal rounds with success probability $1/2$
have no-success probability $1/16$. If all four outcomes equal one shared
fair Bernoulli variable, each round still has marginal success probability
$1/2$, but the chance of no success is $1/2$. Repeated versions of the same
unsuccessful hypothesis can resemble the latter dependence pattern. The
authored example does not diagnose a measured COAT dependence structure; it
shows why marginal responsiveness alone cannot specify a multi-round success
probability.~\citep{exec19_coat_reference}

A public prototype should predeclare the proposal budget, the meaning of a
useful factor and the evidence supplied as feedback. If a quantitative
convergence claim is sought, it must state whether successes are independent
with a fixed probability or satisfy an explicit conditional lower bound given
the preceding history. It also needs a justified contraction statement and an
exact or controlled probability bound. Without those premises, a maximum-round
stop is a resource limit and an unchanged blanket is a search observation.
Neither is upgraded to ``all causal variables recovered.'' Separate final
tasks should remain outside the feedback loop so that repeated improvement on
the development sample is not counted again as fresh validation.

\subsection{Graph preservation and effect estimation retain their conditions}
\label{exec19:coat-graph}
The paper's graph extensions add further structure beyond a single target's
blanket. Its Proposition 8 states adjacency and arrowhead preservation within
a subset under blanket inclusion in the enlarged variable set and the stated
maximally informative PAG setup. The adjustment development additionally uses
amenability, forbidden-variable and blocking conditions. Proposition 15
preserves an adjustment set under its specified blanket-inclusion premise.
These are conditional results about the constructed measured-variable graph;
they are not general certificates for arbitrary hidden-confounder or
cross-world questions. A proposed adapter must preserve uncertain marks and
refuse an effect query when the applicable premises or adjustment route are
unresolved.~\citep{exec19_coat}

The marketing evaluation uses Hillstrom and Retail Hero in a confounded
evaluation setting with synthesized social-media text. For each experiment,
the authors choose the largest valid adjustment sets from the generated PAG
and average estimated effects when there are multiple such sets; experiments
are repeated three times. The extracted Table 7 reports COAT RMSEs of
0.47 and 0.50 for DeepSeek-V3, and 1.90 and 0.78 for GPT-4.1-mini on the
two datasets. The human-factor N-IPW reference reports 1.32 and 1.39.
Thus the comparative result depends on model and comparator: the first
GPT-4.1-mini number does not beat that human-factor reference. These are
source-reported outcomes in the paper's design, not locally reproduced
estimation or confidence-interval calibration.~\citep{exec19_coat}

For a proposed lending study, average error, graph validity, adjustment
availability and interval coverage should be recorded separately. A failure
to recover a valid set must remain in the request denominator. Averaging
effects over several selected sets does not remove the need to justify their
validity or account for selection and dependence. A generated-text benchmark
can expose parser and graph failures at low cost, but it does not establish
that authentic customer narratives measure the same factors.

\subsection{Financial examples motivate a clock and cost audit}
\label{exec19:coat-finance}
COAT's financial case uses Microsoft price and news observations from
2006--2009, with the first 200 trading days for elicitation and the next 400
for evaluation. The target describes the sign of a following-week return.
The reported strategy fits a rolling regression and takes a long or short
position according to its predicted sign. This is an observational trading
illustration. Its reported performance does not identify the effect of a bank
action on a customer outcome. The printed rolling-window summation limits
also need clarification against the executable source; this proposal does not
silently repair them and assume the code did the same thing.~\citep{exec19_coat}

A faithful public reproduction would bind news-publication times, data
availability, revision policy, normalization windows, forecast horizon and
trade execution time. Overlapping weekly outcomes require a declared
dependence treatment. It would retain transaction costs, financing and
failed-query costs before interpreting a trading return as a useful economic
comparison. This is an inexpensive way to test the proposed workflow's
meaning before requesting bank data, rather than treating a published
historical return as bank ROI.

The proposed integration is therefore a staged representation-assistant
boundary: factor instructions and source passages enter a reviewed ontology;
the parser produces values with provenance and explicit missingness; graphical
analysis preserves uncertainty and scientific refusals; downstream estimation
records its own assumptions. The paper's project link is a source-code lead,
not an audited repository, license or runtime. Native artifact review,
measurement-error evaluation and finite stopping guarantees remain separate
work. Public or stipulated gold factors can demonstrate whether this route
merits those later experiments, while independent acceptance and authentic
bank truth remain later evidence.~\citep{exec19_coat,exec19_coat_reference}


% SOURCE docs/release/2026-10-29/source/docs/writing-repair/2026-10-29/src/c008a9ffbc3b9168.tex CHAR 0
\section{Simulated customers and causal assistants need different experiments}
\label{exec20:pa-scope}

A bank considering a new servicing message may ask a language model two
different questions. How would a customer with a given history respond to the
message? Which features of the history justify a causal explanation of the
response? The first question invites a simulated respondent; the second asks
for assistance with a causal model. A convincing answer to either question
does not establish the other. This distinction determines what a public-data
prototype should test before an expensive bank experiment is justified.

Simmons and Hare review language models as subpopulation representative
models: systems intended to reproduce behavior or opinions of a specified
population. They organize prompting, fine-tuning, grounding data and deployment,
and distinguish direct comparison with observed responses, human assessment
and qualitative aggregate evaluation. Their five guiding criteria are fidelity,
necessity, robustness, sensitivity and fairness. In this use, a demographic
description is a proposed conditioning variable; it is not evidence that the
generated respondent represents a customer selected by the bank's actual
contact and response processes.~\citep{exec20_population}

The selected May 2026 revision of Cai, Liu and Song instead manipulates inputs
to a causal assistant. It compares variable names, numerical information and
their removal, then studies instruction fine-tuning on pairwise discovery
tasks. These experiments concern the assistant's behavior under controlled
input conditions. They can guide a prototype evaluation without turning an
assistant's answer into customer-level intervention evidence.
\citep{exec20_attribution}

\subsection{A precise simulation can retain a population error}
\label{exec20:pa-population}
For a fixed customer segment, response definition and historical period, let
$p$ be the true probability of a specified binary response. Let $q$ be the
probability that a fixed language-model configuration generates that response.
Both probabilities refer to the same question, but to different populations
of outcomes. Suppose, for this project derivation, that $M$ independent
decoder draws have common success probability $q$, and let $\widehat q_M$
be their mean. Then
\begin{equation}
 E\{(\widehat q_M-p)^2\}
   =(q-p)^2+\frac{q(1-q)}{M}.
 \label{exec20:pa-mse}
\end{equation}
This follows by writing $\widehat q_M-p=(\widehat q_M-q)+(q-p)$;
the cross term has expectation zero. More draws reduce Monte Carlo error
around the model's own response law. They do not remove the discrepancy
between that law and the target population. If $p$ must itself be estimated
from a survey, uncertainty and selection error in that estimate are additional
obligations.~\citep{exec20_population_attribution_reference}

For an exact assumed example, take $p=1/2$, $q=3/5$ and $M=10{,}000$.
The squared discrepancy is $1/100$, while the simulation variance is only
$3/125{,}000$. If a specified response earns USD 100 and each contact costs
USD 55, the model suggests USD 5 per contact while the stipulated population
law gives USD $-5$. This is a generated planning example, not an estimated
bank return. Moreover, a response forecast is not automatically an incremental
effect of contacting a customer: that decision also needs the relevant
no-contact potential outcome and identification assumptions.
\citep{exec20_population_attribution_reference}

The survey's distinction between consistency, pattern correspondence and
diversity is useful here. Producing a syntactically valid answer or a plausible
customer narrative is a weaker test than matching joint response patterns and
within-segment variation. Its distinction between apparent and actual model
sensitivity also cautions against treating a demographic prompt as a complete
description of the information on which the model relies.
\citep{exec20_population}

A proposal-stage comparison can therefore use a frozen public response
dataset, with its population, dates, response construct and reuse terms
documented. The comparison should include simple empirical or statistical
response models, preserve subgroup denominators, and evaluate held-out
responses and their joint variation. That establishes a bounded descriptive
task. A second, separately declared experiment must test an intervention
question under a known generated mechanism or a defensible public causal
design. Simulated respondents are useful for testing interfaces, eliciting
possible failure cases and comparing designs; their apparent realism supplies
no additional observed bank customers.

\subsection{Input ablations identify a controlled assistant experiment}
\label{exec20:pa-attribution}
In the attribution paper, context and model conditions are held fixed while
names associated with domain knowledge and numerical data are varied. Write
$a_{kd}$ for the correctness probability under a fully specified input
construction, where $k=1$ retains meaningful variable names, $d=1$ retains
numerical input, and zero denotes the corresponding removal. The probability
law must state whether it averages over tasks, decoder randomness, or both.
The paper's conditional and marginal attributions can then be written as
\begin{align}
 \mathrm{CAK}&=a_{11}-a_{01}, &
 \mathrm{CAD}&=a_{11}-a_{10}, \nonumber\\
 \mathrm{MAK}&=a_{10}-a_{00}, &
 \mathrm{MAD}&=a_{01}-a_{00}.
 \label{exec20:pa-factorial}
\end{align}
Here CAK and CAD retain the other input, while MAK and MAD compare against
the no-name, no-data baseline. The selected paper implements removal of
name-associated cues through arbitrary replacement names and removal of
data through prompts containing names alone. Its prompting discussion also
distinguishes wording that asks an assistant to ``suggest'' a relationship
from wording that may induce noncommittal responses.
\citep{exec20_attribution}

These are concrete changes to an interface. Replacing names does not delete
the model's learned parameters, remove every semantic cue in the context,
or prove that all other aspects of difficulty are unchanged. The defensible
experimental target is the effect of the declared prompt transformation on
assistant performance. Calling that transformation knowledge removal requires
additional construct-validity evidence. A refusal must remain an observed
outcome; a more encouraging prompt should not improve a score merely by
removing refusals from its denominator.

The equality $\mathrm{MAD}-\mathrm{MAK}=\mathrm{CAD}-\mathrm{CAK}$
follows algebraically from the four cells. It does not independently test
whether the prompt transformations isolate knowledge and data. To see what
the conditional comparisons count, define the project-derived interaction
\begin{equation}
 I=a_{11}-a_{10}-a_{01}+a_{00},\qquad
 \mathrm{CAK}+\mathrm{CAD}=(a_{11}-a_{00})+I.
 \label{exec20:pa-interaction}
\end{equation}
For assumed cell probabilities $(a_{11},a_{10},a_{01},a_{00})=
(9/10,3/5,2/5,3/10)$, the total gain is $3/5$, the interaction is $1/5$,
and the two conditional attributions sum to $4/5$. Adding them as independent
portions of improvement counts the interaction twice. The ordered paths
$\mathrm{CAK}+\mathrm{MAD}$ and $\mathrm{CAD}+\mathrm{MAK}$ each
equal the total gain. These exact values illustrate the estimand; they are
not measurements from the paper.~\citep{exec20_population_attribution_reference}

For a banking assistant, a paired factorial design should freeze the model,
system instructions, permitted tools, decoding settings and response parser.
Each generating mechanism should appear under all declared input conditions,
with randomized ordering and separate conversation state. The test should
record correct, incorrect, ambiguous and refused outputs. Names should include
both informative and deliberately misleading descriptions, so that a familiar
story cannot substitute for the numerical assumptions that identify a graph.

\subsection{Residual independence must survive the instruction boundary}
\label{exec20:pa-residual}
The fine-tuning section supplies names, definitions and numerical results from
discovery algorithms in the instruction. It describes four pairwise scenarios
through residual correlation, including an undefined direction and the two
possible directed relations. It then teaches an assistant to use knowledge
when the numerical direction is undefined. This description raises a specific
interface question: does the actual numerical artifact test independence under
the required model, or only ordinary zero correlation? The inspected paper
does not settle the native implementation of that boundary.
\citep{exec20_attribution}

An exact project example shows why that distinction matters. Let independent
$X,U$ each take values $-1$ and $1$ with equal probability, and set $Y=2X+U$.
The forward population least-squares residual is $U$, independent of $X$.
The reverse slope and residual are
\begin{equation}
 b=\frac{\operatorname{Cov}(X,Y)}{\operatorname{Var}(Y)}=\frac25,
 \qquad R=X-bY=\frac{X-2U}{5}.
 \label{exec20:pa-reverse}
\end{equation}
Both $\operatorname{Cov}(X,U)$ and $\operatorname{Cov}(Y,R)$ are zero.
The latter is the ordinary least-squares orthogonality condition, not an
independence result. Direct enumeration gives
\begin{equation}
 E(R^2Y^2)=\frac9{25},\quad E(R^2)E(Y^2)=1,
 \quad \operatorname{Cov}(R^2,Y^2)=-\frac{16}{25}.
 \label{exec20:pa-dependence}
\end{equation}
Thus the reverse residual is dependent on its regressor despite zero linear
correlation. A correlation-only instruction can lose the relevant distinction.
This is not a claim that LiNGAM or the authors' executable artifact implements
that faulty rule; neither executable was run for this comparison.
\citep{exec20_population_attribution_reference}

The proposed adapter should consequently carry the numerical method, its
assumptions, the actual dependence test, sample size, uncertainty and refusal
state into the assistant's input. If the engine reports an unidentified
direction, a language-model suggestion can become a separately labeled
hypothesis or request for additional evidence. It must not silently turn that
state into an identified edge. A benchmark that rewards confident completion
of an undefined scenario needs a different acceptance interpretation from a
banking assistant that preserves uncertainty.

\subsection{Reported performance and a public qualification design}
\label{exec20:pa-design}
The attribution study uses nine datasets with fifteen column permutations
each, giving 135 task instances. Its TDR divides correctly predicted pairs by
true causal pairs; it also reports false discovery rate, structural Hamming
distance and F1. Those measures answer different questions, so high TDR alone
cannot certify that the assistant avoids additional false edges. The reverse
task describes keeping numerical values intact while changing the proposed
ordering or adjacency. Such a transformation is a useful input challenge,
but it does not itself generate observations from a new intervention law.
\citep{exec20_attribution}

For instruction fine-tuning, the paper reports 300 pairs producing 1,200
training instances and 106 pairs producing 424 testing instances. It names
Mistral-7B-v0.2 with a LoRA adapter in its training description. Table 4 reports
overall accuracy 0.90 for the fine-tuned model, compared with 0.82 for
GPT-4 turbo, while the undefined scenario scores 0.58 versus 0.66. The
comparison therefore supports a scoped reported gain, not superiority on
every scenario. Its later table uses ``Mixtral'' wording despite the main
training description; exact weight and code identity require resolution
before adoption. Equal scores on three additional tasks are not a universal
no-forgetting result. The paper itself leaves wider benchmark transfer,
finite-sample behavior, missing data and uncertainty-weighted evidence
combination as limitations or future work.~\citep{exec20_attribution}

Repeated prompts quantify one kind of instability, not new independent
domains. Under an illustrative law in which all fifteen outcomes within each
dataset are identical and the nine datasets are independent, 135 calls contain
nine independent units. This extreme is an authored example, not an estimate
of the paper's dependence. A proposed comparison should preserve dataset and
mechanism groups, hold out causal pairs and generating mechanisms, and report
both within-task variability and variation across tasks.
\citep{exec20_population_attribution_reference}

The first prototype can be inexpensive. Its generated mechanisms should
include known directions, observational ambiguity, latent-confounding
violations, and the residual witness above. Its baselines should include a
numerical engine alone, a deterministic translation of that engine's result,
an untuned assistant and the proposed tuned assistant. Success requires
faithful preservation of numerical evidence and uncertainty, alongside any
improvement in explanation or discovery; agreement with a convenient answer
template alone is insufficient. Execution traces should separate the numerical
result, its transformation into an instruction and the assistant's final
claim, making a failed boundary localizable.

Public population-response evaluation and generated causal tasks establish
different merits. Together they can justify whether further prototype work is
worth funding, while keeping actual bank response truth, institutional review
and investment acceptance as later decisions. Neither a large synthetic
respondent pool nor a high assistant score replaces the banking question's
population, information clock, causal assumptions and uncertainty contract.


% SOURCE docs/release/2026-10-29/source/docs/writing-repair/2026-10-29/src/7839d2c7f619ac3f.tex CHAR 0
\section{Measuring whether a causal assistant has improved}
\label{exec20:validity-scope}

A bank comparing two causal assistants needs to distinguish a better answer
from an answer that is easier for its grader to approve. An assistant may
produce a polished diagram, an executable adjustment formula and a confident
explanation while changing the population, conditioning on a variable measured
after the decision, or overlooking an unmeasured common cause. A test that
rewards the format can improve while the scientific error remains. Conversely,
an assistant that correctly refuses a nonidentified query can look worse under
a rubric that rewards every numerical response. The evaluation must therefore
specify the scientific service, the observation used to assess it and the
conditions under which the observation measures that service.

Lin's validity-guided workflow separates an LLM used as a research tool, an
evaluation target, a human simulator and a cognitive model. Its six stages
connect the intended claim to instrument development, experimental design,
documented execution, dependence-aware analysis and bounded interpretation.
The source is a methodological tutorial, with illustrative supplementary code;
it does not report an original empirical validation of a banking assistant.
For this proposal, its most useful contribution is the discipline of validating
the measurement before interpreting a model comparison. A classifier requires
evidence that it performs its specified task; fluent explanations do not
establish a psychological mechanism, and simulated customer answers require
separate evidence of correspondence to the intended population.
\citep{exec20_validity}

\subsection{The target is correctness under a declared information contract}
\label{exec20:validity-target}

Consider a retention-offer question whose request identifies the eligible
customers, the offer time, the balance horizon and the information available
before the offer. The assistant must preserve these elements when it proposes
a graph and an estimand. A useful benchmark includes several legitimate outputs:
a supported identification formula, an explicit effect set, or a refusal that
names the missing assumption or observation. Correctness is judged against the
stated graph class and information, rather than against a single preferred
diagram. If the scenario does not determine an edge direction, selecting the
examiner's favorite direction is not the target.

Let $A\in\{0,1\}$ denote the assigned assistant workflow, with the base model,
retrieval access and available tools specified for each arm. Let $C=1$ mean that
an answer meets a frozen semantic criterion. Let $G=1$ mean that a grader marks
the answer as passing. These are distinct binary variables. The criterion
could require a correct adjustment set and an explicit warning about an
unidentified query; other criteria, such as factual grounding and valid code,
should also be reported individually. A single binary endpoint is introduced
here only to expose a measurement problem. It cannot capture every consequence
of a wrong scientific answer.

Define $p_a=\Pr(C=1\mid A=a)$, grader sensitivity
$s_a=\Pr(G=1\mid C=1,A=a)$ and grader specificity
$c_a=\Pr(G=0\mid C=0,A=a)$. All probabilities refer to the same declared
scenario and answer population. By conditioning on correctness,
\begin{equation}
 q_a:=\Pr(G=1\mid A=a)
 =s_a p_a+(1-c_a)(1-p_a)
 =(s_a+c_a-1)p_a+1-c_a .
 \label{exec20:validity-measurement}
\end{equation}
This is an authored measurement model, not an empirical estimate from Lin's
tutorial. Random assignment of $A$, consistency of the assigned workflows and
an appropriate no-interference condition can identify a causal effect on the
observed grade for the experimental population. Those conditions alone do not
make $q_1-q_0$ an effect on semantic correctness. The relation between $G$ and
$C$ must also be supported. \citep{exec20_validity_reference}

\subsection{A better grade can conceal an unchanged answer}
\label{exec20:validity-example}

Suppose true correctness equals $p_0=p_1=3/5$. Impose grader properties
$(s_0,c_0)=(4/5,9/10)$ for the first workflow and
$(s_1,c_1)=(19/20,4/5)$ for the second. The more polished output is easier to
recognize when correct but also more likely to be approved when wrong. These
properties are assumptions of the example, not observations about a particular
grader. Equation~\eqref{exec20:validity-measurement} gives
\begin{equation}
 q_0=\frac{13}{25}=0.52,\qquad
 q_1=\frac{13}{20}=0.65,\qquad
 q_1-q_0=\frac{13}{100},\quad p_1-p_0=0.
 \label{exec20:validity-illusion}
\end{equation}
More samples can estimate the observed 13-percentage-point improvement more
precisely while leaving its interpretation wrong. Randomizing the workflow
does not remove this error because the workflow itself changes the relation
between the answer and its measurement.

If the two calibration quantities are known for each arm and
$s_a+c_a\ne1$, inversion gives
\begin{equation}
 p_a=\frac{q_a+c_a-1}{s_a+c_a-1}.
 \label{exec20:validity-inverse}
\end{equation}
It recovers $3/5$ in both arms of the example. Even an invariant grader need
not preserve the size of an improvement. With $s=4/5,c=9/10$, changing true
correctness from $1/2$ to $7/10$ changes observed passing by $7/50$; dividing
by $s+c-1=7/10$ recovers the true change of $1/5$. When $s+c=1$, the grade
contains no information about $p$: with $s=c=1/2$, true correctness $1/10$
and $9/10$ both produce a pass probability of $1/2$.
\citep{exec20_validity_reference}

Calibration is therefore an information requirement. In practice, $s_a$ and
$c_a$ must be estimated using held-out adjudication under the intended arms,
scenario population and rubric. Both correct and incorrect cases need support;
verification only of accepted answers cannot identify both quantities without
additional sampling information. Uncertainty in the three inputs must propagate
through the ratio in~\eqref{exec20:validity-inverse}. A denominator close to
zero makes correction unstable. Known bounds on calibration can instead be
propagated to a set of compatible correctness rates, intersected with $[0,1]$;
an unsupported point correction should not be substituted. Model updates,
subgroup differences and changes in answer style can invalidate transported
calibration. A finite example with known probabilities proves none of those
transport assumptions.

The economic consequence is immediate but remains conditional. Suppose, solely
for planning, that an additional correct result is worth USD~100 and the new
workflow costs USD~8 more per request. Using true correctness, the incremental
expected value is $100(p_1-p_0)-8=-8$ dollars. Substituting grades produces
$100(q_1-q_0)-8=5$ dollars. The sign reverses without any improvement in the
scientific service. Real costs of wrong advice, useful refusal and delayed
decisions may differ sharply; the example supplies neither a bank price nor an
investment recommendation. \citep{exec20_validity_reference}

\subsection{Repeated answers and the population of inference}
\label{exec20:validity-dependence}

Lin distinguishes individual observations, experimental units and units of
analysis, and recommends pairing answers to shared items and respecting grouped
observations. This guidance must be applied to the actual design. A fixed model
version is not automatically a randomly sampled model from a population, and a
prompt condition is not automatically the smallest randomized unit. If a
workflow is randomized across independent scenarios, the scenario can be the
assignment unit. If both workflows are run on each scenario with controlled
order and fresh contexts, the paired difference can be the analysis unit. These
designs target different sources of variation from repeatedly sampling decoder
outputs for one fixed question. \citep{exec20_validity}

Aggregation preserves the dependence that remains between groups. To make this
qualification explicit, consider an authored real-valued score model
\begin{equation}
 Z_{jr}=U+V_j+e_{jr},\qquad
 \overline Z_j=U+V_j+\frac{1}{R}\sum_{r=1}^R e_{jr},
 \label{exec20:validity-shock}
\end{equation}
where all distinct components are independent, have zero mean, and have finite
variances $\sigma_U^2,\sigma_V^2,\sigma_e^2$. Here $j$ indexes groups and $r$
repeated scores; the common $U$ can represent a shared evaluation-environment
shock. This is a separate real-valued example, not the binary measurement model.
For distinct groups $j,k$,
\begin{equation}
 \operatorname{Var}(\overline Z_j)
 =\sigma_U^2+\sigma_V^2+\frac{\sigma_e^2}{R},\qquad
 \operatorname{Cov}(\overline Z_j,\overline Z_k)=\sigma_U^2.
 \label{exec20:validity-averages}
\end{equation}
With variances $1/4,1/2,1$ and $R=20$, group means have variance $4/5$,
covariance $1/4$ and correlation $5/16$. Averaging has reduced decoder noise
but has not made the groups independent. \citep{exec20_validity_reference}

Accordingly, the tutorial's shorthand about aggregation eliminating correlation
cannot justify independent group means without assumptions about the shared
environment and scenario sampling. No fixed cluster-count threshold guarantees
calibration, and balanced group sizes alone do not establish a universal
equivalence between aggregation and every mixed model. For a fixed list of
models, descriptive paired differences over a declared scenario collection may
be the honest result. Generalization to unseen scenarios requires a sampling or
transport argument; generalization to future model versions requires another.

\subsection{A public evaluation design for the proposed service}
\label{exec20:validity-design}

The proposed first evaluation uses generated causal mechanisms with known
identification status, together with public scenario descriptions whose
permitted uses and limitations are documented. Generated mechanisms allow an
exact answer key for graph compatibility and intervention queries. Public
business descriptions test translation and missing-information detection, but
their prose does not supply hidden causal truth. The two sources should retain
separate outcome denominators and support different claims.

During development, analysts can revise prompts, retrieval rules and the
semantic rubric on a designated subset. The final evaluation then freezes those
choices and withholds whole mechanisms and scenario families, including their
paraphrases and graph-isomorphic variants. This extends the tutorial's
development/held-out distinction to leakage risks specific to causal reasoning.
Formatting perturbations, variable renaming and column order are diagnostic
conditions within a scenario family; they are not new independent domains.
\citep{exec20_validity}

Before running the comparison, specify whether the intervention changes only
the prompt, or also the model, retrieval corpus, numerical engine and repair
budget. A paired design can compare two frozen workflows on the same held-out
scenarios while controlling order and preventing shared conversation state.
The prespecified outcomes should include semantic correctness, supported
refusals, unsupported assertions, parsing failures, computation failures and
coverage of requests. Report all assigned requests, with a justified treatment
of missing or invalid outputs, so that a workflow cannot improve its headline
score merely by discarding difficult cases.

The grading calibration sample should span both arms and include independently
adjudicated correct and incorrect answers under a prespecified verification
sampling scheme. Blind adjudicators to workflow identity where feasible, retain
disagreements and report both the original grade and the corrected or bounded
semantic result. If expert judgment is itself fallible, disagreement and
verification cannot simply be renamed truth. Known generated answers support
some calibration tasks; interpretation of an ambiguous business request may
require a set of admissible answers.

Finally, preserve the exact requests, outputs, parser decisions, model and tool
versions, retrieval snapshots, collection times, token counts and retries.
Report grading and adjudication costs alongside execution cost and decision
quality. A deterministic checker for a formal formula, independent expert
adjudication and an LLM grader are substantive alternatives with different
coverage and error properties. This design specifies what a prototype should
demonstrate using public and generated evidence before an expensive bank trial;
it does not report a completed experiment or validated bank economics.



% SOURCE docs/release/2026-10-29/source/docs/writing-repair/2026-10-29/src/2ab3aa8b6f11881a.tex CHAR 0
\section{Causal reasoning checklists need matched errors and honest denominators}
\label{exec21:check-scope}
A bank may ask an agent to turn a servicing problem into a causal diagram, identify missing assumptions and explain whether an intervention is justified. A correct final answer on one familiar question is weak evidence for that workflow: the model might recognize the story while missing the mechanism. CausalityCheck, by Li, Duo, Gao and Su, offers a useful response by testing several tasks and perturbations around a seed question. The inspected original is the 27-page July 2026 ACL Findings article. Its contribution is a textual evaluation framework; it does not establish an identified banking effect or a qualified diagram-construction service.~\citep{exec21_causalitycheck}

\subsection{What the checklist tests, and what generates its gold answers}
\label{exec21:check-design}
The framework crosses four tasks---original questions, answerability judgment, process judgment and causal-validity assessment---with four settings: original problem solving, rewritten understanding, irrelevant disturbances and virtual-vocabulary reasoning. Answerability tasks remove key information; process tasks alter a proposed chain; causal-validity tasks distinguish four categories combining correctness and reliability. Virtual reasoning replaces real entities with fictional names. This design is attractive for testing whether a system attends to the provided context and whether its answer survives intended meaning-preserving changes.~\citep{exec21_causalitycheck}

The source of the labels matters. The paper uses GPT-5 to expand CausalProbe seeds and reports two rounds of manual checks with ten trained graduate students. One seed produces 36 instances across the 16 task-by-setting groups because the tasks have different numbers of answer variants. The reported sequence is 1,200 seeds, 43,200 generated instances, a selected 25\% subset of 10,800, and 8,920 retained instances after filtering. The printed cell counts reproduce that final total. These are the authors' reported construction and project-checked arithmetic, not a reproduction of the original model calls or manual judgments.~\citep{exec21_causalitycheck,exec21_causalitycheck_reference}

The four task totals in the printed figure are 482, 2,052, 2,190 and 4,196. Consequently, equal weighting of tasks or cells differs from weighting all retained instances equally. With cell accuracies $r_j$ and cell sizes $n_j$, the two summaries are
\begin{equation}
 R_{\rm macro}=\frac{1}{J}\sum_{j=1}^J r_j,
 \qquad
 R_{\rm micro}=\frac{\sum_{j=1}^J n_j r_j}{\sum_{j=1}^J n_j}.
 \label{exec21:check-weights}
\end{equation}
The paper describes its table entries as averages of related units. A product evaluation must specify the unit and aggregation weights, retain seed identifiers, and avoid treating multiple transformations of the same seed as independent observations. A reported model ranking is otherwise inseparable from the chosen mixture of tasks. Recovering raw instances and their grouping is a prerequisite for a new statistical comparison.~\citep{exec21_causalitycheck,exec21_causalitycheck_reference}

The paper's annotation table also supports a more qualified conclusion than blanket excellent agreement. It prints answerability kappa values of 0.72 and 0.78 in two settings and causal-validity values from 0.59 to 0.65. Several other entries exceed 0.8. Agreement is informative about labeling consistency, but does not make a causal interpretation true. For a bank, gold labels should distinguish claims entailed by a stipulated mechanism from domain judgments needing an owner or additional evidence. Disputed reliability categories should remain visible rather than be converted silently into certain causal facts.~\citep{exec21_causalitycheck}

\subsection{A difference of accuracies is not a frequency of reasoning errors}
\label{exec21:check-paired}
The article uses the difference between answerability and process-judgment accuracies to diagnose causal-chain errors. Table 1 reports 82.27\% and 59.46\% respectively for one model, producing a 22.81 percentage-point difference. A useful diagnostic contrast is not automatically the probability that a model gets answerability right but the chain wrong. To see the distinction, consider an ideal matched population of problems and let $A$ denote a correct answerability judgment and $B$ a correct process judgment. Then
\begin{equation}
 P(A)-P(B)=P(A\cap B^c)-P(A^c\cap B).
 \label{exec21:check-difference}
\end{equation}
The subtraction nets two different kinds of discordance. Without the joint outcomes, it does not identify either one.~\citep{exec21_causalitycheck,exec21_causalitycheck_reference}

For marginal probabilities $a=P(A)$ and $b=P(B)$, the compatible matched-error range is
\begin{equation}
 \max(0,a-b)\leq P(A\cap B^c)\leq\min(a,1-b).
 \label{exec21:check-bounds}
\end{equation}
Using $a=0.8227$ and $b=0.5946$ as an authored matched-design illustration gives a range from 22.81\% to 40.54\%. At the lower endpoint, the joint probabilities for both-correct, $A$-only, $B$-only and neither are $0.5946,0.2281,0,0.1773$. At the upper endpoint they are $0.4173,0.4054,0.1773,0$. Both reproduce the same marginals. The exact reference checks nonnegativity and every marginal. Because the paper's task populations and counts differ, these illustrative bounds are not asserted to be empirical bounds for its actual model responses. Matching must be established first.~\citep{exec21_causalitycheck_reference}

A negative accuracy difference is equally important. With matched marginals 71.68\% and 76.80\%, their difference is $-5.12$ points, but $P(A\cap B^c)$ can still range from zero to 23.2\%. The article's negative challenge score therefore does not by itself demonstrate an absence of chain errors. A proper evaluation should retain paired predictions and gold judgments at seed level and report both discordance directions, uncertainty under the seed sampling law and the intended causal level of each task. Even a well-measured process-label error rate is not a rate of incorrect interventional conclusions unless the labels establish that connection.~\citep{exec21_causalitycheck,exec21_causalitycheck_reference}

\subsection{Filtering changes the decision population}
\label{exec21:check-selection}
The paper removes incomplete fields, ambiguous relationships and sensitive material, and reports 1,880 exclusions from the 10,800 selected instances. Such filtering may be appropriate for a clearly defined benchmark, but the resulting accuracy concerns accepted instances. A production agent will also encounter underspecified, disputed and unanswerable requests. Its refusal policy, escalation burden and error rate on those requests matter to the bank. A benchmark should therefore publish the sampling and exclusion flow, reasons by category and the scope of the reported denominator. The original article's filter is not evidence that the excluded requests are solved.~\citep{exec21_causalitycheck,exec21_causalitycheck_reference}

An authored illustration makes the economic issue explicit. Start with 100 requests, retain 50 and answer 45 of those correctly. Selected accuracy is 90\%. If nothing is known about the excluded half, full-population accuracy lies between 45\% and 95\%. Suppose a correct action yields USD 10, an incorrect action loses USD 30, and every attempted action costs USD 1. At success probability $p$, stipulated expected margin is
\begin{equation}
 V(p)=10p-30(1-p)-1=40p-31.
 \label{exec21:check-value}
\end{equation}
Applying selected accuracy to every request gives USD 5, whereas the compatible full-population margins range from USD -13 to USD 7. These are sensitivity bounds under a deliberately incomplete observation law, not a confidence interval or a measured bank return. If the system abstains on excluded cases, the evaluation must instead model that action's cost and residual human workload explicitly.~\citep{exec21_causalitycheck_reference}

The same lesson applies to source filtering and data-generation quality. GPT-generated items followed by manual review may be useful, but a competing model outperforming the generator does not prove that generator-specific artifacts are absent. The paper's Appendix G offers such a relative-performance argument. A stronger design would compare independently generated and human-authored items, blind the adjudicators to model identity, retain difficult disagreements, and separate development from evaluation seeds. A generator can bias every model while still failing to rank first.~\citep{exec21_causalitycheck}

\subsection{Entity replacement is a stress test, not a unique causal diagnosis}
\label{exec21:check-rename}
The article reports higher performance on virtual-vocabulary tasks and discusses a case where replacing a familiar entity changes a model's handling of a negative conclusion. This motivates a practical perturbation test. It does not uniquely identify memorization or any single internal failure mechanism: replacement can also change tokenization, semantic cues, question difficulty and response-policy triggers. Repeated identical answers to one prompt support repeatability under that setting, not the exclusion of those alternative explanations.~\citep{exec21_causalitycheck}

For a banking prototype, the controlled object should be a known mechanism with explicit variables, action clock and admissible observations. Multiple names and paraphrases can be assigned to the same mechanism while keeping its causal query fixed. A second axis can change the mechanism while retaining familiar surface vocabulary. Contrasting these axes tests whether the system follows the stipulated structure rather than a remembered story. Labels should include answerable, unidentified, unsupported and internally inconsistent cases. This is a proposed experiment motivated by the paper, not an experiment that the library has already run.

\subsection{How to use the framework in the proposal}
\label{exec21:check-integration}
CausalityCheck's task matrix is a candidate design component for evaluating business-question translation, assumption elicitation and diagram explanation. Its textual gold labels should not become an oracle for graph identification, counterfactual effects or hidden-confounder exclusion. The article itself acknowledges limited treatment of structural rigor and counterfactual reasoning. The proposal should combine checklist perturbations with exact finite causal models, known nonidentification pairs and independently adjudicated domain tasks. That combination tests both linguistic robustness and whether the requested effect follows from the available information.~\citep{exec21_causalitycheck}

Before adopting the package or dataset, recover a pinned repository and data revision, trace the transformation and filtering code, inspect license and redistribution terms, and bind item-level provenance. The current review has inspected the genuine article, not executed its code or live models. The library, CLI and MCP should expose seed and variant identity, requested causal level, gold-label origin, inclusion status, paired predictions and uncertainty scope. Agents may organize and explain that evidence; a single aggregate accuracy must not authorize a banking intervention. Public mechanism tests can establish whether a prototype merits an expensive institutional trial while preserving the limits of textual benchmark evidence.


% SOURCE docs/release/2026-10-29/source/docs/writing-repair/2026-10-29/src/6512ea5a9df0c021.tex CHAR 0
\section{From failure documents to bowties: useful extraction without causal overclaiming}
\label{exec21:ace-scope}
A bank's operational-risk team may want to turn incident reports and control inventories into a diagram of causes, a critical event, preventive barriers, consequences and mitigating barriers. Such a representation can organize evidence and expose missing controls. It does not yet specify the effect of installing a control. ACE, by Venkatesh, Oudmans and Zurfluh, supplies a useful industrial analogue: three pipelines generate bowtie diagrams from Failure Mode and Effects Analysis documents. The 3 July 2026 publisher edition compares retrieval, OCR and a vision-assisted pipeline, with substantially better results on selected structured tables than on narrative text. Its central contribution for this proposal is a concrete extraction workflow and a set of failure modes to test before treating a diagram as a causal model.~\citep{exec21_ace}

\subsection{What is extracted, and from which evidence?}
\label{exec21:ace-pipeline}
In the retrieval pipeline, Markdown FMEA content is chunked, embedded with all-MiniLM-L6-v2 and indexed with FAISS; part-name queries select context for a structured prompt. The OCR pipeline uses PaddleOCR and Img2Table to convert table images to Markdown before generation. In the vision pipeline, Mistral-Small produces a Markdown table and description for a second language model to convert into bowtie JSON. The reported implementation uses Python 3.9, sentence-transformers and local GGUF inference through llama-cpp-python. These are named dependencies and reported components, not a pinned, independently reproduced installation in this project.~\citep{exec21_ace}

The evaluation first screens eight model variants on four causal-loop extraction texts and ten false-premise questions. It then selects three instruction-tuned model families and builds tabular references for a sensor, valve assembly and shaft from the US Navy reliability handbook. A domain expert constructs those reference diagrams. Narrative evaluation uses dynamic-seal text and three reference bowties. The paper explicitly notes that expert diagrams sometimes contain implicit variables absent from the source. This is a serious labeling distinction: agreement with an expert's augmented interpretation and faithful extraction of explicit text are different tasks.~\citep{exec21_ace}

A bank-facing interface should therefore retain a source span or table-cell identity for each proposed role and edge. It should separately label an explicit statement, an expert-supplied hypothesis, a model inference and an unresolved conflict. If a narrative does not state a mitigating barrier or consequence, the result can be incomplete. A schema requiring every field to contain a plausible phrase would reward invention. ACE identifies precisely this problem in fallback prompts that ask for a mechanism and otherwise insert a generic mechanism label. Its narrative example also lacks enough information to complete the consequence side of the expert bowtie.~\citep{exec21_ace}

The article's attribution license is recorded with the genuine PDF. It does not establish reuse rights for handbook data, model weights, pipeline code or dependencies. Recovering the actual prompts, preprocessing, model hashes, generation settings and evaluation files is a separate step before adopting an implementation. The proposal can adopt the representation and provenance requirements now while keeping executable reuse a conditional choice.

\subsection{Read metrics with their alignment and denominator}
\label{exec21:ace-metrics}
ACE matches predicted and reference labels with SBERT cosine similarity at an empirically selected threshold of 0.75. It reports node and edge precision, recall and F1 by role or edge type. Its reported graph edit distance is a count of false-positive and false-negative nodes and edges after matching, rather than the minimum edit path in the usual graph-theoretic problem. Threshold choice, matching constraints and reference roles are consequently part of the measurement contract. Semantically similar labels cannot be allowed to erase a reversal of causal direction or the distinction between a cause and a barrier.~\citep{exec21_ace}

Table 4's selected configurations range from node F1 0.81 to 0.94 and edge F1 0.58 to 0.93. The paper describes a missed critical event yielding node F1 0.80 but edge F1 0.44 for a shaft case. On narratives, Table 5's average accuracy ratings range from 1.20 to 3.00 on a five-point scale. The strongest clarity rating, 3.80, is not evidence that the same output has correct causal structure. Expert-level records, rating units, missing ratings and aggregation rules remain necessary to interpret the reported means. These are small, selected industrial evaluations, not banking transfer results.~\citep{exec21_ace}

The counts also deserve direct reconciliation. For each of two prompting conditions, hallucinated, unclear and correct-rejection responses are respectively $(21,10,49)$ and $(9,11,60)$, each summing to 80. Hallucinations decline from $21/80=26.25\%$ to $9/80=11.25\%$, a 15 percentage-point reduction. Unclear responses increase by one, despite prose describing them as reduced. Separately, 12 of 32 schema outputs are invalid JSON, or 37.5\%. These experiments have different denominators and cannot be merged into one accuracy measure. The reference calculation checks the printed arithmetic; no model responses are reproduced here.~\citep{exec21_ace,exec21_ace_reference}

Aggregate F1 can conceal a consequential omission. Consider an authored chain of 101 nodes with 100 directed edges from node zero to node 100. A prediction retains every node and 99 correct edges but omits the bridge from node 49 to node 50. Node F1 is one, and
\begin{equation}
 F_{1,\mathrm{edge}}=\frac{2\,TP}{2\,TP+FP+FN}
 =\frac{198}{199}\simeq0.995.
 \label{exec21:ace-path}
\end{equation}
Nevertheless, the predicted graph contains no directed path from the initiating event to the final consequence. For a bank tracing a payment outage through failed controls to a customer loss, a critical-path or control-role error can outweigh many correctly extracted labels. This example motivates reporting decision-critical role and path errors alongside aggregate graph scores; it does not characterize an observed ACE run.~\citep{exec21_ace_reference}

\subsection{Sensitivity indices do not establish interchangeability}
\label{exec21:ace-sensitivity}
The reported Sobol experiments vary prompt strictness, prompt type, context pipeline and selected decoding settings. They attribute much of the output variation to prompt factors and their interactions. The text interprets a near-zero first-order context index as evidence that the three pipelines are interchangeable under equivalent prompts. That conclusion needs more than a first-order variance decomposition.~\citep{exec21_ace}

For independent design factors $P$ and $C$, let $f(P,C)$ be an evaluation score with positive variance. A first-order context index and a total-order context index are
\begin{equation}
 S_C=\frac{\operatorname{Var}(E[f\mid C])}{\operatorname{Var}(f)},
 \qquad
 S_{T,C}=\frac{E[\operatorname{Var}(f\mid P)]}{\operatorname{Var}(f)}.
 \label{exec21:ace-indices}
\end{equation}
These quantities concern variation under a specified experimental input distribution. To see the distinction directly, stipulate independent uniform binary $P,C$ and $f(P,C)=\mathbf 1\{P=C\}$. Both conditional means given $C$ are $1/2$, while the marginal variance is $1/4$. Thus $S_C=0$ but $S_{T,C}=1$. Changing context at fixed prompt changes the score by either one or minus one: the averaged main effect cancels, although the context matters completely through interaction. Exact enumeration verifies all four cells.~\citep{exec21_ace_reference}

This witness does not show that ACE's actual total-order estimate is one. It establishes why a near-zero first-order estimate alone cannot prove equivalence. Even small estimated total effects over three component tables cannot establish agreement on new formats or rare critical failures. The next experiment should freeze categorical encodings, factor probabilities, semantic matching and decoding controls; retain the full sample matrix and all invalid outputs; and compare pipelines on paired held-out documents with prespecified tolerances for critical-role errors, ordinary extraction accuracy, latency and cost. The precise factor list in the expanded Sobol run, including whether decoding factors remain, should be recovered from code rather than inferred from its abbreviated paragraph.

ACE reports identical JSON in ten runs with different seeds for selected fixed configurations and attributes this to quantization and constrained prompts. Repeated identity establishes observed stability in those trials, not truth or its mechanism. A controlled precision comparison, verified seed propagation and decoding configuration, cache checks and raw token/output records would distinguish competing explanations. Ten repetitions of an erroneous diagram remain the same error. Schema-constrained prompting also should not be described as a hard grammar-constrained decoder without evidence of such an implementation.~\citep{exec21_ace}

\subsection{From a role diagram to an intervention question}
\label{exec21:ace-intervention}
An unsigned arrow is insufficient to price an intervention. In a deliberately minimal example, both structural laws $Y=A$ and $Y=1-A$ have topology $A\to Y$, with binary action $A$ and loss event $Y$. Their intervention contrasts $E[Y\mid do(A=1)]-E[Y\mid do(A=0)]$ are respectively one and minus one. If a loss event costs an assumed USD 10 and the action costs USD 2, the incremental margin is
\begin{equation}
 V=-10\{E[Y\mid do(A=1)]-E[Y\mid do(A=0)]\}-2,
 \label{exec21:ace-value}
\end{equation}
giving USD $-12$ or USD 8. These two laws illustrate missing mechanism information in a common topology; they are not claimed to have identical observational distributions or to satisfy the same signed barrier interpretation. A diagram extractor has not supplied either law merely by naming an edge.~\citep{exec21_ace_reference}

For a bank, the translation layer must additionally identify the action that can be changed, the time at which it is assigned, the eligible unit and population, the loss horizon, the data-generating and observation processes, and any confounding or feedback. A generic barrier node might represent an existing control, a proposed policy, a measured proxy or an expert belief. These possibilities imply different models. Time indexing may be necessary before a feedback description can be represented as an acyclic causal query. Unresolved semantics should remain explicit and should prevent an automatic effect estimate.

\subsection{A bounded agentic workflow and public validation route}
\label{exec21:ace-integration}
ACE proposes future agentic correction and knowledge-graph reasoning, while its reported pipelines are single-pass and single-component. The proposal should keep that boundary visible. A candidate workflow retrieves immutable evidence, proposes typed roles and edges, checks schema and source support, and asks a bounded correction stage to address a specific omission or contradiction. It preserves all attempts and unresolved results. A graph database can store those objects, but neither storage nor repeated agent agreement validates a causal claim.~\citep{exec21_ace}

The CLI or MCP contract should expose source hashes and spans, role types, temporal scope, explicit versus inferred status, schema/model/prompt versions, costs and distinct failure states. Invalid JSON, missing source evidence, a disputed expert addition and a scientifically unidentified intervention are different outcomes. No pipeline should convert all four into a generic failure, fabricate a completed bowtie, or silently grant an intervention-analysis capability.

Public documents and generated mechanisms can justify the next experiment before costly institutional evidence is requested. Split by document family or component, not repeated prompts on the same case. Include explicit and implicit causes, absent barriers, contradictory statements, multi-line tables and misleading nearby components. Freeze the reference construction and separate text-supported truth from expert augmentation. Report invalid outputs, critical-role/path loss, false additions, disagreement, correction attempts and total review cost. Add known-mechanism tasks only where the desired intervention is defined, and retain refusal when extraction cannot support it. Such a prototype can establish whether assisted diagram construction merits a bank trial without claiming that the trial's outcome is already known.~\citep{exec21_ace,exec21_ace_reference}


% SOURCE docs/release/2026-10-29/source/docs/writing-repair/2026-10-29/src/375ddbfcde90051c.tex CHAR 0
\section{Causal programs of thought: computation after the question is justified}
\label{exec21:prophets-scope}
A servicing manager asks whether offering a payment holiday today would reduce
losses over the next ninety days. A useful assistant must determine what the offer
changes, whose losses are being compared, and what information is available before
the offer. Only then can it choose a calculation. Shrestha, Malberg and Groh's
\emph{From Causal Parrots to Causal Prophets?} examines a narrower but relevant
part of this workflow: answering formal causal questions whose numerical facts
and causal structure are encoded in a synthetic benchmark. Its causal tree and
program of thoughts provide concrete alternatives to a single verbal reasoning
chain. The reported improvement motivates a public prototype of assisted
calculation. It does not establish that the assistant can recover the right
banking question or identify an effect from available customer records.
\citep{exec21_prophets}

\subsection{Preserve the question before selecting a solver}
\label{exec21:prophets-question}
Let $A$ denote the offer made at the decision time, $Y(a)$ a default indicator
over the specified horizon under offer $a$, $L(a)$ the monetary loss if default
occurs, and $C(a)$ the cost of that offer. A possible decision target is the
incremental expected cost
\[
 \Delta=E[C(1)-C(0)]+E[Y(1)L(1)-Y(0)L(0)].
\]
This is a proposed banking target, not an estimate supplied by the paper. It
requires a declared eligible population, treatment versions, horizon and
accounting convention. A probability contrast alone is insufficient if the
offer also changes exposure or recovery. A question about customers who accept
an offer differs from a question about making the offer: acceptance can be a
post-offer variable. An assistant that silently changes the action can produce
correct arithmetic for a different decision.

The paper samples 1,000 commonsensical questions from CLADDER and excludes
backdoor-adjustment questions because those questions do not require the
formal calculations and multistep reasoning being studied. It generates
anti-commonsensical and nonsensical variants using GPT-4o and reports manual
checks of those variants. The former replace an outcome with an unrelated
noun; the latter replace variables with generated four-letter strings. These
transformations probe dependence on familiar associations. They do not test
whether a free-form bank request supplies enough information to define an
intervention, an observation process or an eligible population.
\citep{exec21_prophets}

That exclusion matters to the proposed architecture. Selecting a defensible
adjustment set is precisely one of the tasks between a diagram and a numerical
effect estimate. A score on the retained tasks cannot be used as evidence that
this omitted step works. The evaluation should therefore distinguish a
question with a supplied structural model and probabilities from one requiring
model elicitation, identification or data assessment. The former supports a
calculation test. The latter may properly end with an unresolved assumption or
a request for additional evidence.

\subsection{Alternative thoughts are proposals, not independent evidence}
\label{exec21:prophets-search}
The study evaluates thirteen model variants, with the reported experiments
using temperature 1.0. Its causal tree of thoughts follows a six-step causal
reasoning prompt. At selected steps, it generates three alternatives, asks the
model to score them from one to ten, and continues with the highest-scoring
alternative. Other steps have a single continuation. The retained state includes
the question and preceding selected thoughts. This is a greedy selection
procedure, rather than an enumeration of all compatible causal models or a
confidence set over graphs.~\citep{exec21_prophets}

For the servicing example, candidate diagrams might disagree about whether
financial distress precedes the offer, whether acceptance mediates its effect,
or whether a collections contact is observed after treatment. Model scores can
rank these proposals for examination. They cannot establish the timing or
direction of an edge. The interface should retain the source statement,
time index and unresolved alternatives when an edge matters to identification.
A graph rejected by the model's own score can still express a scientifically
relevant alternative. Premature pruning would hide that uncertainty from the
downstream solver.

Self-consistency is a different procedure: the paper samples several reasoning
chains and votes over their answers. The errors of chains generated by the same
model from the same question can be shared. Unanimity measures agreement among
those outputs, not independent confirmation of the causal assumptions. A
comparison of verbal reasoning, voting and tree search should consequently
measure selection mistakes and unresolved alternatives as well as final-answer
accuracy. Their additional model calls and reviewer effort must enter the same
cost accounting as code-assisted methods.~\citep{exec21_prophets}

\subsection{Generated code computes a supplied interpretation}
\label{exec21:prophets-calculation}
The causal program of thoughts asks a model to generate DoWhy/Python code for
the question, executes the code through a Python REPL, and returns the computed
estimate to the model for its final answer. The illustrated workflow generates
synthetic data using numerical facts in the question. Its attraction is clear:
a numerical library can perform operations that a verbal chain might calculate
incorrectly. The causal interpretation still enters through the question,
chosen graph, generated data and selected library operation.
\citep{exec21_prophets}

Suppose a question supplies probabilities under a declared structural law.
Artificial rows sampled from those probabilities may approximate a functional
of that law. They are not newly observed customers and do not strengthen its
empirical support. Increasing the generated sample size can reduce a Monte
Carlo approximation error while leaving the uncertainty about the law intact.
Where the supplied finite probability table permits direct evaluation, an exact
formula checker is therefore a substantive comparator. It avoids introducing
sampling variation solely to recover a quantity already specified numerically.
This is a proposed comparison, not a replication of the study's code.

The same boundary applies to selecting an estimator. A successful function call
does not demonstrate exchangeability, positivity, correct measurement or an
appropriate counterfactual model. The paper's program error analysis separates
incorrect extracted graphs, incorrect library function calls and code execution
failures. Those distinctions should survive in the product result. A tool error
can invite a bounded retry; an unidentified target requires new assumptions,
information or a refusal to provide the requested effect. Treating both as
syntax errors encourages an agent to search until some numerical output appears.
\citep{exec21_prophets}

\subsection{Read gains with their task and denominator}
\label{exec21:prophets-scores}
Table 4 reports GPT-4o overall accuracies of 82.0\% for the baseline,
86.4\% for the causal tree and 90.8\% for the causal program. These are
overall averages across the three variants. The corresponding commonsensical
subset averages are 84.3\%, 87.5\% and 92.5\%. Keeping the two columns distinct
prevents a favorable overall number from being assigned to the wrong
population. Subtraction of the rounded overall scores gives gains of 4.4 and
8.8 percentage points. The causal program's counterfactual scores are 92.5\%,
90.2\% and 87.4\% across commonsensical, anti-commonsensical and nonsensical
questions, compared with baseline scores of 70.1\%, 68.8\% and 64.1\%.
The corresponding differences are 22.4, 21.4 and 23.3 percentage points.
\citep{exec21_prophets,exec21_prophets_reference}

Those differences are descriptive calculations over published rounded entries.
The printed tables alone do not supply the paired outcomes needed for a statistical comparison. Section~\ref{exec21:prophetsrepo-results} subsequently examines saved question-level answers and reproduces the selected averages; the sampling law and reference validity remain separate requirements. The variants derive
from common underlying questions, so they should not automatically be treated
as independent observations. Nor is an error selected for explanation a random
draw from all deployment errors. The appropriate sampling unit, model revision,
prompt, retries and treatment of failed outputs must accompany a comparison.

Several appendix counts need reconciliation before secondary rates are computed.
Table 5's sampled query-type counts sum to 999, whereas its rung counts sum to
1,000; both corresponding original-dataset totals are 8,690. The response-pattern
counts in Tables 7 and 8 sum to 995 and 999 for ten-chain and five-chain voting.
Table 10 lists error counts of 10, 37 and 13, totaling 60, although its preceding
description states an analysis of 50 samples. Overlapping categories, missing
outputs or transcription errors are possible explanations, but the inspected
text does not settle them. Recovering identifiers and raw records is necessary
before choosing a denominator or asserting an error-category rate. The subsequent repository review resolves the one-item query-type discrepancy from the saved identifiers, but leaves the vote histories and error-category annotations unresolved.
\citep{exec21_prophets,exec21_prophets_reference}

The paper also reports unsuccessful memorization probes and acknowledges that
GPT-4o generated the perturbations on which it was evaluated. Failed recall of
three questions does not establish that related training data were absent.
The authors explicitly restrict their study to formal causal reasoning and one
synthetic benchmark, and call for evaluation on further tasks. These limitations
make the paper useful evidence for a narrowly specified prototype. They prevent
its scores from establishing general causal discovery, faithful diagram
construction or business-language translation.~\citep{exec21_prophets}

\subsection{Audit the reference calculation as well as the agent}
\label{exec21:prophets-arithmetic}
Appendix D.3 illustrates why a benchmark answer cannot be accepted solely because
it is printed beside a worked solution. Write $m_{xv}$ for the numerical outcome
probability displayed for treatment $x$ and mediator $v$, and $\pi_x$ for the
displayed probability that the mediator is one. The appendix supplies
$m_{00}=0.28$, $m_{01}=0.60$, $m_{10}=0.35$, $m_{11}=0.74$,
$\pi_0=0.57$ and $\pi_1=0.22$. Evaluating its step-4 expression literally gives
\begin{align}
 \sum_{v=0}^1 m_{0v}
 \{P(V=v\mid X=1)-P(V=v\mid X=0)\}
 &=(m_{01}-m_{00})(\pi_1-\pi_0)\notag\\
 &=(0.60-0.28)(0.22-0.57)=-0.112.
 \label{exec21:prophets-literal}
\end{align}
Its step-6 expansion is a different expression. Literal evaluation gives
\begin{align}
 &\{0.22(0.60-0.28)+0.78(0.74-0.35)\}\notag\\
 &\quad-\{0.57(0.60-0.28)+0.43(0.74-0.35)\}
 =0.0245,
 \label{exec21:prophets-expanded}
\end{align}
not the printed 0.0173. These calculations use exact rational arithmetic;
they reproduce the displayed inputs without executing any downloaded program.
\citep{exec21_prophets,exec21_prophets_reference}

The numerical discrepancy is only part of the issue. The narrative supplies
death probabilities but names the outcome as survival, and its displayed
counterfactual notation must be reconciled with the intended natural indirect
effect. Equations~\eqref{exec21:prophets-literal}
and~\eqref{exec21:prophets-expanded} consequently identify inconsistencies;
neither is declared the corrected benchmark truth. Section~\ref{exec21:prophetsrepo-reference} later reconciles the step-4 numerical functional with the saved target while retaining the outcome-coding and counterfactual ambiguity. Establishing scientific truth still requires the intended coding, structural model, counterfactual target and actual generator. A library can faithfully calculate the wrong expression or
use the right expression with reversed outcome coding. A benchmark protocol
must admit a disputed-reference outcome and preserve it outside the accuracy
denominator until adjudicated under a predeclared rule.

In the banking example, a similar reversal between default and survival could
change the action recommendation. Correctness must therefore be attached to
the complete target and its units, rather than merely to a signed decimal.
Independent arithmetic checks are valuable precisely because they inspect a
different failure mechanism from language-model agreement. They still do not
supply the scientific assumptions that make a mediation or intervention
functional appropriate.

\subsection{A public prototype with separate tests for each claim}
\label{exec21:prophets-validation}
The proposal can investigate these methods without first requesting an expensive
bank experiment. Begin with public, known-law cases whose variables, decision
times, mechanisms and targets are fixed before model evaluation. Split by
underlying mechanism and document family, so paraphrases and renamed variables
do not leak between development and evaluation. Retain three related tasks:
extract a question and its evidence; determine whether a supplied model supports
the requested target; and calculate that target when it is supported. Include
cases with missing assumptions, contradictory coding and several compatible
diagrams, for which a single effect is not the correct output.

Compare a no-tool answer, a direct formula checker, calls to a constrained
causal library and a generated-code agent on the same frozen cases. Review
source-span fidelity, temporal roles, missing variables, graph alternatives,
identification/refusal, numerical error and final decision separately. Count
invalid code, timeouts, retries and abstentions in declared denominators. A
correct final yes/no answer must not conceal an unsupported graph or a wrong
estimand that happened to yield the same sign. Blinded reference adjudication
should precede scoring, and unresolved reference disagreements should be visible.

The comparison also needs model and tool revisions, seeds or decoding settings,
raw responses, generated programs, dependencies and all failed attempts.
Question-level paired results support uncertainty calculations that respect the
shared cases. Token use, tool calls, latency and human correction time belong
in a full workflow comparison. Useful merit means meeting a predeclared
combination of decision quality, justified refusal and total cost; it is not
established by the paper's accuracy improvement alone. These are proposed
public qualification steps, with no live-model result claimed here.

\subsection{A proposed CLI and MCP boundary}
\label{exec21:prophets-contract}
A request should distinguish a calculation from a proposed interpretation. It
should carry the decision time, eligible population, intervention, outcome
coding, horizon, source identities and spans, graph alternatives, estimand and
explicit assumptions. A calculation request additionally needs its actual
probability table or dataset and a declared sampling law. The response should
return the selected target, unresolved conditions, calculation provenance and a
typed disposition. Missing scientific information, conflicting evidence,
unsupported identification, invalid code and numerical failure are different
outcomes requiring different next actions.

For a generated-code prototype, the proposed execution boundary permits only
declared inputs and bounded computation in a restricted environment. Generated
code, dependency identities and tool inputs must be inspectable, and network
access or policy actions must not be implicit consequences of a calculation.
A formula checker or typed library call can supply a simpler comparator with
a smaller execution surface. The paper motivates evaluating these alternatives;
its use of a REPL is not evidence that this proposed interface has been
implemented or qualified. The following repository review inspects selected prompts, saved evaluations and source rights, identifying both reusable records and defects that rule out unchanged executable adoption. Complete artifacts, dependencies and a qualified runtime remain prerequisites for a new comparison.
\citep{exec21_prophets}

The resulting investment question is concrete: can constrained calculation
support improve a faithfully specified causal workflow on public cases at an
acceptable total cost? A positive result would justify a more expensive
prototype or institutional study while preserving the additional assumptions
that authentic banking decisions require.


% SOURCE docs/release/2026-10-29/source/docs/writing-repair/2026-10-29/src/f325b5256237efd2.tex CHAR 0
\section{What the saved causal-reasoning records establish}
\label{exec21:prophetsrepo-scope}
The public repository linked by Shrestha and colleagues permits a more precise
assessment of the preceding paper-level questions. Its selected 23 March 2025
commit contains prompts, inference helpers, notebooks and result tables.
Twelve selected files have been checked against their pinned Git blob identities.
The following source and saved-result findings supersede the preceding section's
unresolved availability statements for those selected materials. They resolve
some reporting questions while revealing why a saved benchmark score and a
reproducible, qualified agent are different claims. No upstream program or
model was executed in this inspection.~\citep{exec21_prophets_repo}

\subsection{Reproduce the score before interpreting the gain}
\label{exec21:prophetsrepo-results}
Each of the three RQ3 result tables has 1,000 unique question identifiers. The
identifiers, reference labels, query types and rungs align across variants.
Direct comparison of the eight method-answer columns with their stored labels
reproduces all twenty-four subset averages in Table 4. In particular, baseline,
tree and program correct counts are 843, 875 and 925 on commonsensical cases;
824, 865 and 903 on anti-commonsensical cases; and 793, 853 and 896 on nonsense
cases. Across the 24,000 selected method-by-variant final-answer entries, three
zero-shot entries are blank: two in the anti-commonsensical table and one
in the nonsensical table. All other entries are yes or no. The blank entries
remain in their 1,000-question denominators and count as not correct. This is a replay
of scoring over saved outputs, not a new model evaluation or independent
validation of the reference labels.~\citep{exec21_prophets_repo,exec21_prophets_repo_reference}

The raw sample includes 138 treatment-on-the-treated questions, resolving the
one-item discrepancy with the paper's printed count of 137. Its rung counts
are 405 association, 197 intervention and 398 counterfactual questions. The
paired records also show the errors hidden by an average gain. On commonsensical
cases the program corrects 141 baseline errors but loses 59 answers the baseline
got right. Those counts are 152 and 73 on anti-commonsensical cases, and 182 and
79 on nonsense cases. Net gains are therefore 82, 79 and 103 correct answers
per thousand, consistent with the reported subset averages.
\citep{exec21_prophets_repo_reference}

For a proposed banking assistant, the regression cases deserve their own
review. A method that improves average accuracy can still fail on a particular
decision category or on cases where the simpler method succeeds. The records
permit paired analysis by question. They do not supply a sampling justification
for treating all paraphrased variants as independent deployment observations.
Uncertainty and decision-cost comparisons should respect shared question and
mechanism identities, and should not interpret repeated reference labels as
independent adjudications.

Some appendix discrepancies remain unresolved. The five-chain and ten-chain
self-consistency reasoning fields are empty in all three saved tables. Final
majority labels cannot reconstruct individual chain votes or explain the
reported vote-pattern totals of 999 and 995. The selected records likewise do
not contain the error-category annotations needed to reconcile sixty listed
errors with a fifty-case description. These limitations concern particular
missing records; they do not negate the twenty-four reproduced accuracy
averages.~\citep{exec21_prophets_repo,exec21_prophets_repo_reference}

\subsection{A numerical target can reconcile while its meaning remains disputed}
\label{exec21:prophetsrepo-reference}
Question 2883 matches the penguin example in Appendix D.3. Its stored numerical
target is approximately $-0.11217602178843022$. Applying the displayed
step-4 functional to the record's unrounded probabilities yields
$-0.11217602178843015$, agreeing to within $10^{-15}$. The rounded calculation
of $-0.112$ in the preceding section is consequently consistent with the saved
numerical target. The positive expansion printed in the paper is not a
reproduction of that target.~\citep{exec21_prophets_repo_reference}

This resolves the numerical comparison, not the semantic problem. The stored
question still describes death probabilities while asking about survival, and
the stored formal expression still needs reconciliation with the intended
natural indirect effect. Even the saved reasoning string contains an arithmetic
expression inconsistent with its stated result. A scorer can faithfully compare
an agent's answer with the stored ``yes'' while leaving these problems untouched.
The original generator, outcome coding, polarity and intended intervention
must be inspected before relabeling the case or calling its reference
independently correct.~\citep{exec21_prophets_repo}

In a financial application, analogous disagreement could concern whether a
field denotes default, repayment, positive recovery or a negative loss. A small
numerical tolerance is useful after those semantics are fixed. It cannot settle
them. The proposed result contract should retain both the recorded numerical
functional and the unresolved interpretation, rather than silently converting
agreement with a stored number into a validated decision recommendation.

\subsection{The example generator and the displayed graph disagree}
\label{exec21:prophetsrepo-generator}
The program prompt supplies an illustrative drug-assignment example. Read
literally as an intended generator, it draws assignment $A$ with probability
one half, then independently draws adherence $T$ and low cholesterol $Y$
conditional on $A$. Its probabilities are
\[
 \begin{aligned}
 P(T=1\mid A=0)&=0.79,& P(T=1\mid A=1)&=0.43,\\
 P(Y=1\mid A=0)&=0.54,& P(Y=1\mid A=1)&=0.52.
 \end{aligned}
\]
The generated outcome does not depend on adherence once assignment is fixed.
Yet the displayed DoWhy graph contains only $T\to Y$ and treats adherence as
the intervention. Assignment, which drives both generated variables, is omitted.
The example also contains explanatory text that is not valid Python inside
its purported code illustration. These are properties of the selected prompt
text; the prompt module itself is a Python string, and no failure rate for
generated programs follows from inspecting it.~\citep{exec21_prophets_repo}

An exact project calculation isolates the scientific issue independently of
syntax. Under the independent-noise generator just described,
\begin{equation}
 E[Y\mid do(T=1)]-E[Y\mid do(T=0)]=0,
 \qquad
 E[Y\mid T=1]-E[Y\mid T=0]=\frac{6}{793}.
 \label{exec21:prophetsrepo-confounding}
\end{equation}
The second contrast is about $0.00757$. It follows from mixing over the common
parent $A$, with $P(T=1)=0.61$; the intervention leaves the outcome equation
unchanged. These are exact properties of an authored interpretation of the
displayed generator, not clinical findings, a native DoWhy run or a claim about
every generated answer in the study.~\citep{exec21_prophets_repo_reference}

For the servicing decision, the same failure occurs if an assistant generates
offer acceptance and subsequent loss from distress but passes a graph containing
only acceptance and loss to an estimator. Larger artificial samples would make
the associational calculation more stable without repairing its causal target.
The public prototype should therefore compare the generated mechanism, supplied
graph and requested estimand before permitting estimation. Exact known-law
examples provide an inexpensive way to test this consistency requirement.

\subsection{Reference reuse requires a different execution contract}
\label{exec21:prophetsrepo-reuse}
The selected inference helper has a concrete call mismatch: its self-consistency
function accepts four parameters while its caller supplies five. Within the
helper, the completion call also omits the client required by that routine.
The tree branch in the common runner is unimplemented, with tree logic instead
appearing in a separate notebook. These observations qualify the selected
source revision. They do not establish which revision or interactive notebook
state generated the saved result tables.~\citep{exec21_prophets_repo}

The README recommends installation through a pip requirements command, but the
file is a Conda environment export for macOS on ARM, including Python 3.12,
DoWhy 0.8 and platform-specific build strings. A supported environment cannot
be inferred from that instruction. The notebooks retain fixed-seed sample
construction but also partial example runs, several scoring denominators and
model names that do not bind a provider snapshot. The program runner hands
generated code to a REPL; its use of input sanitization does not establish the
proposed resource, network or tool-permission boundaries.
\citep{exec21_prophets_repo}

The repository carries a complete MIT grant for its author's material. Rights
for inherited benchmark data, dependencies, model access and downstream uses
remain separate. The defensible proposal decision is to retain the prompts
and saved records as research references and reject unchanged executable
adoption. A future implementation should repair call and client plumbing,
use strict typed parsing, record invalid outcomes and retries, bind dependency
and model identities, and separate known-law computation from unsupported
causal interpretation. A constrained library call or exact formula checker
remains a necessary comparator to arbitrary generated code.
\citep{exec21_prophets_repo}

The next public evaluation can now start from authenticated question identifiers
and paired outputs while retaining disputed references and missing histories.
It must still establish its own frozen execution environment, held-out cases,
reference adjudication and full workflow costs. This closes a source-availability
and scoring gap while leaving scientific truth and executable qualification as
explicit, testable obligations.

\section{Choose the comparator that isolates the proposed benefit}

The studies support different alternatives because the tasks are different. For extracting a claim, compare the assistant with structured extraction and source-grounded human review. For computing a supplied functional, compare it with a direct formula checker and a constrained numerical call. For generating factors or simulated responses, compare it with observed measurements and ordinary statistical models. A stronger verbal explanation cannot compensate for a changed intervention, a noisy proxy accepted as a confounder or a synthetic population that misses the decision population.

This comparison also determines the next expenditure. Public mechanisms can reveal a wrong target, an invalid separator, a grading artifact or a failed software handoff before a bank trial. Saved outputs can support paired reanalysis, but their reference labels and execution history retain their own obligations. A prototype merits further investment when its benefit survives these credible alternatives, its unfavorable cases remain visible and its full review and failure costs are counted. The operating service must then preserve those scientific distinctions through the library and its interfaces, which is the subject of the next part.


% SOURCE docs/release/2026-10-29/source/docs/causal_banking_release_candidate_2026-10-29.tex CHAR 28090

\part{Operating the analytical service}
\label{reading:part-practice}
The analyst should be able to follow the same study from source records
to the financial conclusion. A research library, a command-line tool and an
assistant using the Model Context Protocol are alternative ways to conduct that
study. Their design follows from the population, timing, intervention and
uncertainty established in the preceding parts.


